\documentclass[11pt,a4paper]{report}
\usepackage[utf8]{inputenc}
\usepackage[T1]{fontenc}
\usepackage[spanish,es-noshorthands,es-tabla]{babel}
\usepackage{lmodern}
\usepackage[margin=2.2cm]{geometry}
\usepackage{amsmath,amssymb}
\usepackage{siunitx}
\sisetup{output-decimal-marker={.},range-phrase=--,per-mode=symbol}
\usepackage{booktabs,array,tabularx,multirow}
\usepackage{enumitem}
\setlist{nosep,leftmargin=*}
\usepackage[dvipsnames]{xcolor}
\usepackage[most]{tcolorbox}
\usepackage{tikz}
\usetikzlibrary{arrows.meta,positioning,decorations.pathmorphing}
\usepackage{hyperref}
\hypersetup{colorlinks=true,linkcolor=NavyBlue,urlcolor=NavyBlue}
\usepackage{fancyhdr}
\pagestyle{fancy}\fancyhf{}
\fancyhead[L]{\small Física Nuclear Experimental}
\fancyhead[R]{\small\nouppercase{\leftmark}}
\fancyfoot[C]{\thepage}
\usepackage{titlesec}
\titleformat{\chapter}[hang]{\Huge\bfseries\color{NavyBlue}}{\thechapter.}{0.5em}{}
\titlespacing*{\chapter}{0pt}{-20pt}{18pt}

% ---------- Niveles de relevancia ----------
% CLAVE: imprescindible (ideas del resumen oficial, fórmulas que hay que saber)
\newtcolorbox{clave}[1][]{enhanced,breakable,colback=red!4,colframe=BrickRed,
  fonttitle=\bfseries,title={$\bigstar$ Clave\ifx&#1&\else: #1\fi},sharp corners=south}
% FÓRMULA destacada
\newtcolorbox{formula}[1][]{enhanced,colback=NavyBlue!4,colframe=NavyBlue,
  fonttitle=\bfseries\small,title={#1},halign=center}
% DATO: valores de referencia, no memorizar
\newtcolorbox{dato}[1][]{enhanced,breakable,colback=gray!6,colframe=gray!50,
  fontupper=\small,fonttitle=\bfseries\small,coltitle=black,
  title={\textsf{i}\ Dato\ifx&#1&\else: #1\fi}}
% EJEMPLO
\newtcolorbox{ejemplo}[1][]{enhanced,breakable,colback=ForestGreen!4,colframe=ForestGreen!70!black,
  fontupper=\small,fonttitle=\bfseries\small,title={Ejemplo\ifx&#1&\else: #1\fi}}

\newcommand{\dEdx}{\ensuremath{\mathrm{d}E/\mathrm{d}x}}
\newcommand{\imp}[1]{\textbf{\color{BrickRed}#1}}

\title{\Huge\bfseries Física Nuclear Experimental\\[0.4em]\Large Apuntes de la asignatura\\[0.3em]\large Máster Interuniversitario en Física Nuclear --- Curso 2026-27}
\author{Resumen elaborado a partir de las diapositivas de los Temas 1--12\\y del ``Repaso de Detectores''}
\date{}

\begin{document}
\maketitle

\chapter*{Cómo leer estos apuntes}
\addcontentsline{toc}{chapter}{Cómo leer estos apuntes}
El contenido está jerarquizado en tres niveles de relevancia, usando como guía el
\emph{Repaso de Detectores} (resumen oficial de la asignatura):

\begin{clave}
\textbf{Imprescindible.} Ideas y fórmulas que aparecen en el resumen oficial o que vertebran el tema.
Hay que saberlas, entenderlas y saber aplicarlas.
\end{clave}

\textbf{Texto normal: necesario.} Explicaciones, matices y desarrollo que hay que conocer para
entender lo clave y responder preguntas de razonamiento.

\begin{dato}
\textbf{Dato de referencia.} Tablas, valores numéricos, nombres de instrumentos y detalles
históricos o técnicos. Conviene tener el orden de magnitud, pero no memorizar.
\end{dato}

\begin{ejemplo}
Aplicaciones numéricas para fijar ideas.
\end{ejemplo}

\tableofcontents

\chapter*{Visión general: ideas que hay que tener claras}
\addcontentsline{toc}{chapter}{Visión general: ideas que hay que tener claras}
\markboth{Visión general}{}

\begin{clave}[Cómo interacciona cada tipo de radiación]
\begin{itemize}
\item \textbf{Partículas cargadas}: interaccionan \imp{de forma continua} (miles de colisiones
coulombianas con los electrones del medio). En el caso de $e^-$ y $e^+$, la
\imp{radiación de frenado (bremsstrahlung)} domina sobre la ionización por encima de la
\imp{energía crítica $E_c$}.
\item \textbf{Rayos $\gamma$}: interaccionan mediante \imp{procesos discretos}. Los relevantes son
\imp{efecto fotoeléctrico, dispersión Compton y producción de pares}. En espectroscopía $\gamma$
se usan materiales de \imp{$Z$ alto} para que la sección eficaz fotoeléctrica domine.
\item \textbf{Neutrones}: procesos discretos $\rightarrow$ \imp{reacciones nucleares} (dispersión
elástica/moderación, captura). Su detección es siempre indirecta.
\end{itemize}
\end{clave}

\begin{clave}[Ideas básicas sobre detectores]
\begin{itemize}
\item \textbf{Tipos}: activos/pasivos, directos/indirectos, integradores/contadores,
dosímetros/espectrómetros; por material: gas, semiconductor, centelleador.
\item \textbf{Modos de operación}: pulso y corriente.
\item \textbf{Propiedades generales}: eficiencia, resolución, respuesta temporal (tiempo muerto).
\item \textbf{Detección $\gamma$}: coeficiente de absorción lineal, elección de material, efectos del
tamaño del detector, efectos del material circundante.
\item \textbf{Cadena de detección}: interacción radiación--materia $\rightarrow$ la energía depositada
se convierte en pulso eléctrico $\rightarrow$ amplificación y conformado $\rightarrow$ digitalización
$\rightarrow$ espectro de energía.
\end{itemize}
\end{clave}

\begin{center}
\begin{tikzpicture}[node distance=0.25cm,
  b/.style={draw=NavyBlue,fill=NavyBlue!6,rounded corners,align=center,font=\small,minimum height=1cm,text width=2.3cm,inner sep=2pt}]
\node[b] (a) {Interacción\\rad.--materia};
\node[b,right=of a] (b) {Carga / luz\\$\rightarrow$ pulso};
\node[b,right=of b] (c) {Preamp. +\\conformado};
\node[b,right=of c] (d) {Digitali-\\zación (ADC)};
\node[b,right=of d] (e) {Espectro\\de energía};
\foreach \x/\y in {a/b,b/c,c/d,d/e} \draw[-{Stealth},thick] (\x)--(\y);
\end{tikzpicture}
\end{center}

\section*{Mapa rápido de la asignatura}
\begin{center}\small
\begin{tabularx}{\linewidth}{lX}
\toprule
Tema & Qué hay que llevarse \\
\midrule
1. Interacción & Bethe--Bloch, pico de Bragg, alcance; $e^\pm$: bremsstrahlung, $E_c$, $X_0$; $\gamma$: foto/Compton/pares, $I=I_0e^{-\mu x}$; neutrones: retroceso nuclear.\\
2. Caract. generales & Modos pulso/corriente, espectro diferencial, resolución (FWHM, Fano), eficiencias, tiempo muerto (paralizable/no paralizable), estadística de contaje.\\
3. Gas & Curva de regiones (ionización, proporcional, GM), avalancha, gases, $W\approx 30$ eV.\\
4. Semiconductores & Unión p--n, zona de deplexión $w\propto\sqrt{V}$, $V_{dep}$, Si (cargadas) vs HPGe ($\gamma$), segmentación, telescopios $\Delta E$--$E$.\\
5. Centelleadores & Orgánicos vs inorgánicos, tabla de propiedades, PMT vs SiPM, NaI vs HPGe.\\
6. Neutrones & Reacciones de captura ($^3$He, $^{10}$B, $^6$Li), moderación, ToF, esferas de Bonner, matrices moderadas.\\
7. Electrónica & Preamplificador sensible a carga, conformado CR--RC (ruido, apilamiento), CFD y time walk, SCA/MCA/ADC, digitalización completa.\\
9. Monte Carlo & Muestreo (transformada inversa, $s=-\lambda\ln\xi$), error $\propto1/\sqrt N$, algoritmo de transporte, reducción de varianza, códigos.\\
10. Aceleradores & Electrostáticos y tándem, linac, ciclotrón ($f=qB/2\pi m$), sincrociclotrón, betatrón, sincrotrón; límites de cada uno.\\
11. Vidas medias & Ley de desintegración, $B(\sigma L)$ y unidades Weisskopf, técnicas según escala: electrónicas, Doppler (RDDS, DSAM), fast timing.\\
12. Análisis & SRIM/TRIM, NNDC, ROOT/jsROOT, Gnuplot.\\
\bottomrule
\end{tabularx}
\end{center}

\chapter{Interacción de la radiación con la materia}

\section{Idea general}
Casi todos los detectores funcionan igual: la radiación entra, interacciona con los átomos
(o núcleos) del medio y cede toda o parte de su energía a electrones ligados. Se libera un
número grande de electrones de baja energía (o pares $e$--$h$, o fotones de centelleo) que se
recogen y analizan electrónicamente.

\begin{clave}[Lo que el detector mide siempre son cargas]
En los tres casos (cargadas, fotones, neutrones) lo que el detector mide al final son
electrones/iones, pares $e$--$h$ o fotones de centelleo \emph{liberados por partículas cargadas}.
Fotones y neutrones sólo se ven a través de las partículas cargadas secundarias que producen.
\end{clave}

\begin{center}\small
\begin{tabularx}{\linewidth}{lXXXX}
\toprule
Radiación & Mecanismo dominante & Pérdida de energía & Alcance & Blindaje \\
\midrule
$\alpha$, p, iones & Colisiones coulombianas con $e^-$ (Bethe--Bloch) & Continua; pico de Bragg & Bien definido ($\alpha$ 5 MeV en Si $\approx\SI{25}{\micro m}$) & Papel, $\mu$m--mm de sólido\\
$e^-$, $\beta$ & Colisiones + bremsstrahlung & Continua + radiativa; trayectorias erráticas & Mal definido (mm); $\rho R\approx$ cte & Pocos mm de Al o plástico\\
$\gamma$, X & Fotoeléctrico, Compton, pares & Discreta (sucesos ``catastróficos'') & No definido: $I=I_0e^{-\mu x}$ & Pb, hormigón ($Z$, $\rho$ altos)\\
n & Dispersión elástica, reacciones nucleares & Discreta; retroceso nuclear & No definido; moderación & Material hidrogenado + absorbente (B, Li, Cd)\\
\bottomrule
\end{tabularx}
\end{center}

\begin{dato}[¿Cuánto material hace falta para pararla? (en Si)]
$\alpha$ de 5 MeV: $\approx\SI{25}{\micro m}$; p de 5 MeV: $\approx\SI{0.2}{mm}$; $e^-$ de 1 MeV: $\approx\SI{1.8}{mm}$;
capa hemirreductora de un $\gamma$ de 1 MeV: $\approx\SI{4.7}{cm}$. Cuatro órdenes de magnitud:
un detector de $\alpha$ es una lámina y uno de $\gamma$ un bloque. Para $\beta$ basta 0.1--1 mm;
para $\gamma$ de MeV, 5 cm puede ser insuficiente.
\end{dato}

\paragraph{Elección del material según lo que se quiere medir.}
\begin{itemize}
\item \textbf{Energía}: amplitud proporcional a la energía y \emph{muchos} portadores (pequeñas fluctuaciones estadísticas).
\item \textbf{Tiempo}: recogida rápida de la carga (el número de portadores importa menos).
\item \textbf{Tipo de partícula}: que masa o carga produzcan un efecto distintivo (p.ej.\ $\Delta E$--$E$).
\item \textbf{Ritmo alto}: recuperación rápida (poco tiempo muerto). \textbf{Ritmo bajo}: reducir fondo.
\item \textbf{Trayectoria}: sensibilidad a la posición (segmentación).
\end{itemize}

%=====================================================
\section{Partículas cargadas pesadas}
Todas las cargadas más pesadas que el electrón ($\pi$, $\mu$, p, d, $\alpha$, iones). Interaccionan
con los electrones y núcleos del medio mediante la fuerza de Coulomb, transfiriendo en cada
colisión fracciones mínimas de energía: \textbf{aproximación de frenado continuo (CSDA)}.
Una partícula de 1 MeV sufre $\sim10^5$ interacciones antes de pararse.

\paragraph{Tipos de colisión según el parámetro de impacto $b$ frente al radio atómico $a$:}
\begin{itemize}
\item \textbf{Blandas} ($b\gg a$): excitan o ionizan el átomo entero; pocos eV por colisión; son las más numerosas.
Una fracción $<0.1\%$ puede salir como \emph{luz Cherenkov} si $\beta>1/n$.
\item \textbf{Duras} ($b\sim a$): choque con un electrón concreto que sale con energía considerable
(\textbf{rayo $\delta$}) y deja su propia traza. Son pocas pero se llevan una fracción de energía comparable a las blandas.
\item \textbf{Con el núcleo} ($b\ll a$): dispersión de Rutherford; casi no frena pero desvía (dispersión múltiple).
\end{itemize}

\begin{clave}[¿Quién frena y quién desvía?]
La \textbf{pérdida de energía} la dominan los \textbf{electrones}: hay $Z$ electrones por núcleo y la
energía cedida va como $1/m_{\text{blanco}}$ ($m_e\ll M$). La \textbf{desviación angular} la dominan los
\textbf{núcleos} (dispersión múltiple). Como $M\gg m_e$, la partícula pesada apenas se desvía:
\textbf{trayectoria recta y alcance bien definido}.
\end{clave}

\begin{dato}[Cherenkov]
Umbral $\beta>1/n$; en agua ($n=1.33$) $\beta>0.75$: $\approx0.26$ MeV para $e^-$, $\approx485$ MeV para p.
Ángulo $\cos\theta=1/(n\beta)$, máximo $\approx41^\circ$ en agua. Espectro $\propto1/\lambda^2$ (luz azul).
\end{dato}

\paragraph{Máxima energía transferida a un electrón.}
\[
T_{max}=\frac{2m_ec^2\beta^2\gamma^2}{1+2\gamma m_e/M+(m_e/M)^2}\;\approx\;2m_ec^2\beta^2\gamma^2
\]
Crece como $\beta^2\gamma^2$: con iones lentos los $\delta$ son despreciables; a GeV dominan las fluctuaciones.
En un detector fino los rayos $\delta$ pueden escapar: \emph{energía depositada $<$ energía perdida}.
Para dosimetría se usa el LET restringido (sólo lo depositado cerca de la traza).

\begin{dato}[$T_{max}$]
$\alpha$ 5 MeV: 2.7 keV; p 10 MeV: 22 keV; p 150 MeV: 0.35 MeV; p 1 GeV: 3.3 MeV.
\end{dato}

\subsection{Fórmula de Bethe--Bloch}
\begin{formula}[Poder de frenado (pérdida por colisión)]
\[
-\left\langle\frac{\mathrm{d}E}{\mathrm{d}x}\right\rangle
= K z^2\,\frac{Z}{A}\,\frac{1}{\beta^2}
\left[\frac12\ln\frac{2m_ec^2\beta^2\gamma^2T_{max}}{I^2}-\beta^2-\frac{\delta(\beta\gamma)}{2}\right]
\]
\[ K=4\pi N_A r_e^2m_ec^2=\SI{0.307}{MeV\,cm^2/mol} \]
\end{formula}
\begin{itemize}
\item $z$, $\beta$: carga y velocidad del proyectil. $Z/A\approx0.5$ e $I$ (energía media de
excitación, $I\approx10Z$ eV) caracterizan el medio. $\delta$: efecto densidad.
\item Expresada en $\mathrm{MeV\,cm^2/g}$ (dividiendo por $\rho$) casi no depende del material.
\item Válida para $\beta\gamma\gtrsim0.1$; por debajo: correcciones de capas y de carga efectiva.
\end{itemize}

\begin{clave}[Dependencias de Bethe--Bloch]
\begin{itemize}
\item $\dEdx\propto z^2/\beta^2$: depende de la \textbf{carga y velocidad} del proyectil, \textbf{no de su masa}.
\item A baja energía domina $1/\beta^2$: cuanto más lenta, más ioniza $\Rightarrow$ \textbf{pico de Bragg}.
\item Mínimo hacia $\beta\gamma\approx3$--3.5, casi igual en todos los materiales:
\textbf{partículas de ionización mínima (MIP)}. Después, subida relativista suave (logarítmica), frenada por $\delta$.
\item A igual energía no relativista: $\dEdx\propto Mz^2/E$ $\Rightarrow$ $\Delta E\cdot E\propto Mz^2$:
base de la identificación de partículas con telescopios $\Delta E$--$E$ (Tema 4).
\end{itemize}
\end{clave}

\begin{dato}[MIP en Si]
Una MIP deja $\approx\SI{0.39}{keV/\micro m}$ en Si ($\approx108$ pares $e$--$h$/$\mu$m, $\approx 1.66$ MeV cm$^2$/g).
Protones en Si: mínimo hacia 2--3 GeV. Electrones: mínimo de colisión hacia 1--1.5 MeV.
Valores de $I$: agua 75 eV, aire 85.7 eV, Si 173 eV, Ge 350 eV, Pb 823 eV.
\end{dato}

\paragraph{Correcciones.} Capas (electrones internos no libres a baja $v$); términos de Barkas
($\propto z^3$) y Bloch ($\propto z^4$), distinguen partícula y antipartícula; efecto densidad
(polarización del medio, alta energía). \textbf{Compuestos}: regla de Bragg,
$\frac{1}{\rho}\dEdx=\sum_i w_i\frac{1}{\rho_i}\left(\dEdx\right)_i$ ($w_i$ fracción en masa).

\paragraph{Iones lentos: carga efectiva.} Cuando la velocidad del ion es comparable a la de sus
electrones, captura/pierde electrones y su carga media baja: $z\to z_{ef}$. Esto produce que el
máximo de $\dEdx$ de iones pesados aparezca antes del final. A keV/nucleón aparece el
\textbf{frenado nuclear} (choques elásticos con átomos, daño por desplazamiento). Herramienta: SRIM/TRIM.

\subsection{Alcance y pico de Bragg}
\begin{formula}[Alcance CSDA]
\[ R(E_0)=\int_0^{E_0}\left(\frac{\mathrm{d}E}{\mathrm{d}x}\right)^{-1}\mathrm{d}E \]
\end{formula}
\begin{itemize}
\item Para pesadas, el alcance medio (lo alcanza el 50\%) está muy bien definido y casi coincide con la profundidad alcanzada.
\item \textbf{Escalado}: a igual velocidad, $R\propto M/z^2$ (un d de 10 MeV recorre el doble que un p de 5 MeV).
\item En g/cm$^2$ depende poco del material.
\item \textbf{Pico de Bragg}: la máxima densidad de ionización se da al final del recorrido. Base de la
protonterapia/hadronterapia frente a fotones (dosis máxima en profundidad).
\item \textbf{Straggling}: cada partícula sufre un número distinto de colisiones $\Rightarrow$ alcances
casi gaussianos ($\sigma_R\approx1$--1.2\% de $R$ para p en agua). Iones más pesados: straggling relativo menor ($\propto\sqrt{m_p/M}$), pico más nítido.
\item El alcance disminuye drásticamente con $z$.
\end{itemize}

\begin{dato}[Alcances a 5 MeV en Si]
p: 220 $\mu$m; $\alpha$: 25 $\mu$m; $^{16}$O: 4.3 $\mu$m; $^{40}$Ca: 3.0 $\mu$m; $^{132}$Xe: 2.0 $\mu$m; $^{197}$Au: 1.4 $\mu$m.
Bases de datos: NIST PSTAR/ASTAR/ESTAR y SRIM.
\end{dato}

\begin{ejemplo}[telescopio $\Delta E$--$E$]
Si de 100 $\mu$m delante de un detector grueso. Partículas de 10 MeV: el protón deja
$\Delta E\approx0.85$ MeV ($E_{res}\approx9.15$ MeV), el deuterón $\Delta E\approx1.48$ MeV. A 20 MeV:
0.48 y 0.83 MeV. Cada especie cae en una hipérbola distinta del diagrama $\Delta E$--$E$.
\end{ejemplo}

%=====================================================
\section{Electrones y positrones}
Misma física de colisión que las pesadas pero con diferencias importantes:
\begin{itemize}
\item Suelen ser relativistas (p.ej.\ $\beta^-$).
\item Chocan con partículas de su misma masa: \textbf{grandes desviaciones}, trayectoria tortuosa,
\textbf{alcance difuso} y retrodispersión. Longitud de trayectoria $\gg$ penetración.
\item \textbf{Partículas idénticas}: Bethe se sustituye por Møller ($e^-e^-$) o Bhabha ($e^+e^-$); $T_{max}=T/2$.
\item \textbf{Radiación de frenado (bremsstrahlung)}: la aceleración en el campo nuclear emite fotones.
\end{itemize}

\begin{clave}[Bremsstrahlung y energía crítica]
\[
\left(\frac{\mathrm{d}E}{\mathrm{d}x}\right)_{tot}=\left(\frac{\mathrm{d}E}{\mathrm{d}x}\right)_{col}+\left(\frac{\mathrm{d}E}{\mathrm{d}x}\right)_{rad},
\qquad
\frac{(\mathrm{d}E/\mathrm{d}x)_{rad}}{(\mathrm{d}E/\mathrm{d}x)_{col}}\approx\frac{E\,Z}{700}\ (E\text{ en MeV})
\]
\begin{itemize}
\item La pérdida radiativa crece con $E$ y con $Z^2$ ($\propto Z^2E/m^2$: despreciable para pesadas).
\item \textbf{Energía crítica $E_c$}: aquella en que ambas pérdidas se igualan, $E_c\approx\frac{\SI{800}{MeV}}{Z+1.2}$
(sólidos: $\approx 610/(Z+1.24)$). Por encima de $E_c$ domina la radiación.
\item \textbf{Longitud de radiación $X_0$}: distancia en la que un $e^-$ de alta energía conserva $1/e$ de su energía:
$E(x)=E_0e^{-x/X_0}$. Escala natural de cascadas electromagnéticas y calorímetros; fija también la producción de pares.
\item Para detectar $\beta$: \textbf{materiales de $Z$ bajo} (plástico, Si) y ventanas finas: menos bremsstrahlung y retrodispersión.
\end{itemize}
\end{clave}

\begin{dato}[$X_0$ y $E_c$]
\begin{tabular}{lccc}
\toprule
Material & $X_0$ (g/cm$^2$) & $X_0$ (cm) & $E_c$ (MeV)\\
\midrule
Al & 24.3 & 9.0 & 43\\ Si & 22.1 & 9.5 & 40\\ Ge & 12.6 & 2.4 & 18\\ Pb & 6.3 & 0.56 & 7.3\\
\bottomrule
\end{tabular}\quad
Cociente rad/col a 1 MeV: agua $\approx1\%$, Si $\approx2\%$, Pb $\approx12\%$.\\
Ejemplo: $e^-$ de 1 GeV en 1 cm de Pb ($1.8X_0$) conserva $e^{-1.8}\approx17\%$.
\end{dato}

\paragraph{Alcance de electrones.} No hay alcance único: la curva de transmisión cae poco a poco.
Se usa el \textbf{alcance extrapolado} $R_{ex}$ (prolongar el tramo recto de la curva hasta cortar el eje).
En g/cm$^2$ casi no depende del material: $R_1\rho_1\approx R_2\rho_2$.
\begin{dato}[Katz--Penfold (0.01--2.5 MeV)]
$R_{ex}[\mathrm{mg/cm^2}]=412\,T^{\,1.265-0.0954\ln T}$ ($T$ en MeV). 1 MeV $\to$ 0.41 g/cm$^2$ $\to$ 1.5 mm Al, 1.8 mm Si, 4 mm agua.
\end{dato}

\paragraph{Positrones.} Se frenan como electrones y, casi en reposo, \textbf{se aniquilan}: dos fotones
de \SI{511}{keV} a $\approx180^\circ$ (base de la PET). En un detector $\gamma$ estos fotones pueden
escapar: picos de escape simple y doble.

%=====================================================
\section{Radiación neutra: idea general}
Puede no interaccionar en absoluto o perder su energía en uno o pocos sucesos ``catastróficos''.
No se puede predecir hasta dónde penetrará: se describe por el \textbf{recorrido libre medio}
$\lambda=1/\mu$, no por un alcance.

\section{Radiación gamma}
\begin{clave}[Atenuación exponencial]
\[
I(x)=I_0\,e^{-\mu x},\qquad \mu=\mu_{fe}+\mu_{C}+\mu_{pp}=n\,\sigma_{tot},\qquad x_{1/2}=\frac{\ln2}{\mu},\qquad \lambda=\frac1\mu
\]
$\mu$: coeficiente de atenuación \textbf{lineal} (cm$^{-1}$), suma de los procesos. $\mu/\rho$ (cm$^2$/g)
elimina la dependencia con la densidad; en mezclas se suma ponderado en masa. Válido para
\emph{haz estrecho}; con haz ancho se añade un \textbf{factor de acumulación} (build-up) porque los
fotones dispersados Compton siguen contribuyendo.
\end{clave}

\begin{dato}[$\gamma$ de 1 MeV]
$\mu/\rho\approx0.06$--$0.07$ cm$^2$/g en casi todo (Compton domina): $x_{1/2}$ = 4.7 cm (Si), 2.3 cm (Ge), 0.86 cm (Pb), 9.8 cm (agua).
\end{dato}

\begin{clave}[Los tres procesos]
\begin{center}\small
\begin{tabularx}{\linewidth}{XXX}
\toprule
\textbf{Fotoeléctrico} & \textbf{Compton} & \textbf{Producción de pares}\\
\midrule
Absorbe el fotón entero; toda la energía queda en el detector. &
Choque con un $e^-$ casi libre: el fotón sigue con menos energía. &
$\gamma\to e^-e^+$ en el campo del núcleo; umbral $2m_ec^2=1.022$ MeV.\\
$\sigma\propto Z^{n}/E^{3.5}$, $n\approx4$--5 & $\sigma\propto Z$ (por átomo) & $\sigma\propto Z^2$\\
Domina a baja $E$ y alto $Z$ & Domina en torno a 1 MeV & Domina por encima de unos MeV\\
\bottomrule
\end{tabularx}
\end{center}
Por la fuerte dependencia $Z^{4-5}$, en espectroscopía $\gamma$ se buscan materiales de \textbf{$Z$ alto}
(Ge, NaI, BGO, LaBr$_3$\dots) para maximizar la fracción de fotopico.
\end{clave}

\subsection{Efecto fotoeléctrico}
El fotón es absorbido por un átomo y sale un electrón ligado (preferentemente de la capa K):
\[ T_{e^-}=E_\gamma-E_b. \]
Un electrón libre no puede absorber un fotón (no se conservarían a la vez energía y momento).
La vacante se llena emitiendo rayos X característicos o electrones Auger, que normalmente también
se absorben: \textbf{toda la energía queda en el detector} $\Rightarrow$ fotopico. La sección eficaz
presenta bordes de absorción (K, L\dots).

\subsection{Dispersión Compton}
\begin{formula}[Cinemática Compton]
\[
E'_\gamma=\frac{E_\gamma}{1+\alpha(1-\cos\theta)},\qquad \alpha=\frac{E_\gamma}{m_ec^2},\qquad
T_{e^-}=E_\gamma-E'_\gamma
\]
\[
\text{Borde Compton }(\theta=\pi):\quad T_{max}=E_\gamma\frac{2\alpha}{1+2\alpha}
\]
\[
E'_\gamma(\pi)=\frac{E_\gamma}{1+2\alpha}\xrightarrow{E_\gamma\gg m_ec^2}\frac{m_ec^2}{2}=\SI{256}{keV}
\]
\end{formula}
\begin{itemize}
\item El electrón se lleva una energía continua entre 0 y $T_{max}$: \textbf{continuo Compton} con su \textbf{borde}.
\item El fotón retrodispersado (en el entorno) da el \textbf{pico de retrodispersión} a $E'_\gamma(\pi)$ (entre 184 y 256 keV).
\item Distribución angular: Klein--Nishina; a alta energía, dispersión hacia delante.
\end{itemize}

\subsection{Producción de pares}
Sólo si $E_\gamma>1.022$ MeV, en el campo del núcleo. La energía cinética del par es $E_\gamma-2m_ec^2$.
El positrón se aniquila y emite dos fotones de 511 keV, que pueden escapar del detector:
\begin{itemize}
\item \textbf{Pico de escape simple}: $E_\gamma-511$ keV. \textbf{Pico de escape doble}: $E_\gamma-1022$ keV.
\item Pico de aniquilación a 511 keV (pares creados en el material circundante).
\end{itemize}

\begin{dato}[Fronteras entre procesos]
Fotoeléctrico = Compton: $\approx50$ keV en Al, $\approx560$ keV en Pb. Compton = pares: $\approx15$ MeV en Al, $\approx4.5$ MeV en Pb.
Diagrama $Z$--$E_\gamma$: líneas de igual contribución de procesos vecinos.
\end{dato}

\begin{clave}[Espectro $\gamma$ típico (p.ej.\ $^{137}$Cs en NaI)]
\begin{itemize}
\item \textbf{Fotopico}: toda la energía depositada (en uno o varios pasos).
\item \textbf{Continuo y borde Compton}: un Compton y el fotón escapa (477 keV para 662 keV).
\item \textbf{Valle} entre borde y fotopico: Compton múltiple.
\item \textbf{Pico de retrodispersión} (184 keV para $^{137}$Cs), rayos X del blindaje (Pb), escapes y 511 keV si $E_\gamma>1.022$ MeV.
\end{itemize}
Detector pequeño: escapan secundarios (más Compton y escapes). Detector grande: casi todo acaba en el
fotopico. El material circundante aporta retrodispersión, rayos X y aniquilación.
\end{clave}

%=====================================================
\section{Neutrones (introducción)}
Sin carga: sólo interacciones nucleares (no hay barrera coulombiana), con sección eficaz pequeña.
La energía se deposita a través de los núcleos de retroceso o productos de reacción.
\begin{formula}[Máxima energía transferida en choque elástico frontal]
\[ E_{max}=\frac{4A}{(A+1)^2}E_n \qquad\Rightarrow\quad \text{H: } E_{max}=E_n;\quad \text{C: } 0.28\,E_n \]
\end{formula}
\begin{itemize}
\item El mejor absorbente/moderador es el \textbf{hidrógeno} (centelleadores orgánicos, polietileno). Un bloque de Pb
apenas frena neutrones rápidos ($\sim$2\% por choque).
\item Si la dispersión es isótropa en el CM ($E_n<10$ MeV en H), la energía del protón de retroceso es \textbf{uniforme entre 0 y $E_n$}.
\item Neutrones lentos/térmicos: reacciones de captura $^3$He(n,p), $^{10}$B(n,$\alpha$), $^6$Li(n,$\alpha$) (ver Tema 6).
\end{itemize}

\chapter{Características generales de los detectores}

\section{Clasificación}
\begin{clave}[Tipos de detectores]
\begin{itemize}
\item \textbf{Activos} (lectura electrónica) / \textbf{pasivos} (sin electrónica: placa fotográfica, película radiocrómica, TLD\dots).
\item \textbf{Directos}: la señal sale directamente de la ionización del medio (gas, semiconductor).
\textbf{Indirectos}: la radiación se convierte en luz, que luego se convierte en señal eléctrica (centelleador + fotosensor).
\item \textbf{Integradores} (suman el efecto de muchas interacciones) / \textbf{contadores} (cada partícula individual).
\item Según la información: \textbf{contadores} (número de partículas, p.ej.\ Geiger--Müller),
\textbf{espectrómetros} (distribución en energía, p.ej.\ NaI, HPGe), \textbf{dosímetros} (energía neta depositada).
\item Por material: \textbf{gas, semiconductor, centelleador} (no excluyentes).
\end{itemize}
\end{clave}

\paragraph{Formación de la señal.} Si se deposita suficiente energía se crean portadores:
\[ N=\frac{\Delta E}{W} \]
donde $W$ es la energía media para crear un portador (par $e$--ion, par $e$--$h$, fotón\dots).
\begin{dato}[Escalas de energía por portador]
Gases $\approx30$ eV; semiconductores 3--10 eV (gap 1--5 eV; Si 3.6 eV, Ge 2.96 eV); centelleadores 20--500 eV por fotoelectrón ($\sim10$ eV para estados ópticos);
fonones $\sim$meV. Duración típica de los pulsos: 1 ns -- 10 $\mu$s.
\end{dato}
Cuanto menor es $W$, más portadores y mejor resolución en energía (menor fluctuación estadística relativa).

\paragraph{Materiales según bandas.} Conductores: sin gap, enorme corriente de fondo (inútiles).
Dieléctricos (diamante): gap grande, poca corriente de fuga pero señal pequeña.
\textbf{Semiconductores} (Si, Ge): gap moderado, permiten aplicar campos grandes con corriente de fuga
pequeña; requieren gran pureza.

\section{Modos de operación}
\begin{clave}[Pulso vs corriente]
\begin{itemize}
\item \textbf{Modo pulso}: cada interacción se procesa individualmente. Permite espectroscopía (amplitud $\propto$ energía) y timing.
Problema: \textbf{tiempo muerto}.
\item \textbf{Modo corriente}: se promedian las señales; se pierde la información individual. La corriente
es proporcional a la tasa de dosis. Se usa a tasas muy altas (dosimetría, monitores de haz).
\end{itemize}
\end{clave}

\subsection{Tiempo muerto}
Dos interacciones deben estar separadas un tiempo mínimo $\tau$ para dar señales distintas; si no, la
segunda se pierde o distorsiona la primera (apilamiento). Lo fija el componente más lento de la cadena:
el propio detector en un GM (decenas--cientos de $\mu$s), el ADC en sistemas con multicanal (pocos $\mu$s).

\begin{clave}[Modelos de tiempo muerto ($n$: tasa real, $m$: tasa medida)]
\begin{center}
\begin{tabular}{cc}
\textbf{No paralizable} & \textbf{Paralizable}\\[2pt]
$\displaystyle n=\frac{m}{1-m\tau}$ & $\displaystyle m=n\,e^{-n\tau}$\\[8pt]
Sucesos durante $\tau$ no se cuentan & Sucesos durante $\tau$ no se cuentan\\
ni lo prolongan; $m\to1/\tau$ & \emph{y lo prolongan}; $m\to0$ a tasas altas\\
\end{tabular}
\end{center}
A tasas bajas ($n\tau\ll1$) ambos coinciden: $m\approx n(1-n\tau)$.
En un sistema paralizable $m$ tiene un máximo $m_{max}=1/(e\tau)$ en $n=1/\tau$: una tasa medida corresponde a dos tasas reales.
\end{clave}
\begin{dato}[Ejemplos]
Paralizables: centelleadores con PMT, Si(Li)/Ge(Li) con electrónica lenta, HPGe con preamplificador RC.
No paralizables: contador Geiger, centelleadores con electrónica rápida (PET), rechazadores de apilamiento.
\end{dato}

\section{Espectro de alturas de pulso}
La mayoría de espectrómetros operan en modo pulso: la amplitud es proporcional a la energía
\emph{depositada} (que no siempre es la energía total incidente). Se representa el \textbf{espectro
diferencial} $\mathrm{d}N/\mathrm{d}H$ (cuentas por canal) frente a la amplitud/energía; el número de pulsos entre
$H_1$ y $H_2$ es el área bajo la curva.

\section{Eficiencia}
\begin{clave}[Eficiencia]
\[
\varepsilon_{abs}=\frac{\text{n.º detectadas}}{\text{n.º emitidas}}
=\underbrace{\frac{\text{n.º que llegan}}{\text{n.º emitidas}}}_{\varepsilon_{geo}=\Omega/4\pi}\times
\underbrace{\frac{\text{n.º detectadas}}{\text{n.º que llegan}}}_{\varepsilon_{int}}
\]
\begin{itemize}
\item \textbf{Geométrica}: ángulo sólido. Disco de radio $a$ a distancia $d$ en el eje:
\[ \Omega=2\pi\left(1-\frac{d}{\sqrt{d^2+a^2}}\right),\qquad \Omega\xrightarrow{d\gg a}\frac{\pi a^2}{d^2}\ \text{(inverso del cuadrado)} \]
\item \textbf{Intrínseca}: depende de $E$, $Z$, $\rho$ y espesor. Haz paralelo monoenergético: $\varepsilon_{int}=1-e^{-\mu x}$.
\item Para $\gamma$ se usa la \textbf{eficiencia de fotopico}: área neta del fotopico / $\gamma$ emitidos. Se mide con
fuentes calibradas a varias energías y se \textbf{interpola} con un ajuste (evitar extrapolar).
\end{itemize}
\end{clave}
\begin{ejemplo}
$a=2.5$ cm, $d=10$ cm $\Rightarrow\Omega=0.188$ sr, $\varepsilon_{geo}=1.5\%$. Duplicar $d$ la reduce $\approx4$ veces.
Fuentes extensas o fuera de eje: integración numérica o Monte Carlo.
\end{ejemplo}

\section{Resolución en energía}
Capacidad de distinguir dos energías próximas. La afectan: fluctuaciones estadísticas del número de
portadores (o de fotones de centelleo), colección incompleta de carga/luz y ruido electrónico.
\begin{clave}[Resolución]
\[
\mathrm{FWHM}=2\sqrt{2\ln2}\,\sigma\approx2.355\,\sigma,\qquad R=\frac{\mathrm{FWHM}}{H_0}
\]
Límite estadístico (Poisson con factor de Fano $F$):
\[
\sigma_N=\sqrt{FN},\quad N=\frac{E}{W}\ \Rightarrow\ R_{stat}=2.355\sqrt{\frac{FW}{E}}
\]
$F\approx0.1$--0.2 en gases y semiconductores ($F<1$ porque los sucesos de ionización no son independientes);
$F\approx1$ en centelleadores. Las contribuciones independientes se suman en cuadratura:
$\mathrm{FWHM}^2_{tot}=\mathrm{FWHM}^2_{stat}+\mathrm{FWHM}^2_{ruido}+\mathrm{FWHM}^2_{colección}+\dots$
\end{clave}
Cuando $R$ es pequeña (HPGe) se da directamente la FWHM en keV (p.ej.\ $\approx1.8$ keV a 1332 keV);
en centelleadores se da en \% (NaI $\approx7\%$ a 662 keV). Dos picos se separan si distan más de una FWHM.

\section{Respuesta y resolución temporal}
\begin{itemize}
\item \textbf{Respuesta temporal}: tiempo en formar la señal. Buen ``timing'' = flanco de subida rápido.
\item \textbf{Tiempo de subida} (10\%$\to$90\%), \textbf{tiempo de caída} (90\%$\to$10\%), constante $\tau$ de la caída exponencial.
\item \textbf{Resolución temporal}: anchura de la distribución de diferencias de tiempo entre dos detectores en coincidencia
(CRT, \emph{coincidence resolving time}). Se mide con TAC (analógico) o con marcas de tiempo digitales.
\end{itemize}

\section{Detección de rayos gamma: material y tamaño}
\begin{clave}[Elección del material]
$N=N_0e^{-\mu x}$ y $\mu$ cae rápidamente con la energía. Para $\gamma$ de alta energía hacen falta
materiales de \textbf{$Z$ y densidad altos}.
\begin{itemize}
\item Si a 10 keV: domina el fotoeléctrico ($\mu\approx10^2$ cm$^{-1}$); 300 $\mu$m absorben el 95\%. Bien para X.
\item Si a 1 MeV: casi todo Compton ($\mu\approx0.1$ cm$^{-1}$); en 300 $\mu$m interacciona el 0.3\%; harían falta 3 cm para el 95\% y aun así escapa el fotón dispersado.
\item Por eso: \textbf{Ge} (Z=32) para espectroscopía de alta resolución; \textbf{NaI(Tl)} (I, Z=53), \textbf{BGO} (Bi, Z=83) para alta eficiencia.
\end{itemize}
\end{clave}

\begin{clave}[Efectos del tamaño del detector]
\begin{itemize}
\item \textbf{Muy pequeño}: sólo la primera interacción. Fotopico (fotoeléctrico), continuo Compton y, si $E_\gamma>1.022$ MeV,
\textbf{pico de doble escape} ($E_\gamma-2m_ec^2$): ambos 511 keV escapan.
\item \textbf{Intermedio}: Compton múltiples (cuentas entre borde Compton y fotopico), \textbf{pico de escape simple}
($E_\gamma-m_ec^2$) y doble, y fotopico reforzado.
\item \textbf{Muy grande} (absorción total): idealmente sólo el pico de energía total. Extendido a cascadas: \textbf{espectroscopía
de absorción total} (TAS) para medir alimentaciones $\beta$.
\end{itemize}
\end{clave}

\begin{clave}[Efectos del material circundante]
Rayos X característicos (p.ej.\ Pb, 75--85 keV) por fotoeléctrico en el blindaje; pico de 511 keV por aniquilación de pares
creados fuera; \textbf{pico de retrodispersión} ($\approx$200--250 keV): a $180^\circ$ todas las energías incidentes dan
fotones dispersados con energía similar.
\end{clave}

\section{Etapas de la detección}
\begin{enumerate}
\item Deposición de energía en el medio (para medir energía: alta densidad y $Z$; para trazas: mínima perturbación).
\item Conversión en señal eléctrica (directa o indirecta); carga $\propto$ energía.
\item Amplificación (electrónica o multiplicación interna: avalancha, dínodos). Introduce ruido.
\item \textbf{Conformado}: el pulso estrecho de corriente se transforma en un pulso más ancho (reduce ruido) con máximo
redondeado en el tiempo de pico $T_P$ (facilita medir la amplitud). Si la forma no depende de la señal, amplitud $\propto$ energía.
\item \textbf{Digitalización}: amplitud con ADC; diferencias de tiempo con TDC.
\end{enumerate}

\chapter{Detectores de gas}

\section{Principio de funcionamiento}
Recipiente con gas, dos electrodos aislados y alta tensión entre ellos. La radiación ionizante crea
pares electrón--ion; los iones positivos van al cátodo y los electrones al ánodo. Habitualmente el
cátodo es la pared del recipiente y el ánodo un hilo fino central (geometría coaxial).
\begin{dato}[$W$ (eV/par) en gases]
H$_2$ 37, He 41, N$_2$ 35, O$_2$ 31, aire 35, Ne 36, Ar 26. Valor típico: $W\approx30$ eV.
\end{dato}
Modos: \textbf{corriente} (dosímetro, energía integrada) o \textbf{pulso} (sucesos individuales).

\section{Regímenes de operación}
\begin{clave}[Número de iones recogidos frente a la tensión aplicada]
\begin{center}
\begin{tikzpicture}[x=0.95cm,y=0.42cm,font=\small]
\draw[-{Stealth}] (0,0)--(11,0) node[above left]{tensión $V$};
\draw[-{Stealth}] (0,0)--(0,11) node[above right]{$\log N_{iones}$};
\foreach \x/\l in {0.7/a,2.6/b,5/c,7.3/d,8.9/e} \node at (\x,10.5){(\l)};
\foreach \x in {1.4,3.8,6.4,8.1,9.8} \draw[dashed,gray] (\x,0)--(\x,10);
\fill[green!10] (3.8,0) rectangle (6.4,10);
\draw[thick,BrickRed] (0.1,1.5) .. controls (0.6,4) .. (1.4,4.5) -- (3.8,4.6) -- (6.4,7.8) .. controls (7.3,9.2) .. (8.1,9.5) -- (9.8,9.5) -- (10.5,10.5);
\draw[thick,NavyBlue] (0.1,0.5) .. controls (0.6,2.2) .. (1.4,2.5) -- (3.8,2.6) -- (6.4,5.8) .. controls (7.4,8.5) .. (8.1,9.5);
\node[BrickRed,right] at (4.2,8.7){$\alpha$};
\node[NavyBlue,right] at (5.2,3.9){$\beta$};
\node[align=center,font=\scriptsize] at (2.6,-1.2){cámara de\\ionización};
\node[align=center,font=\scriptsize] at (5.1,-1.2){contador\\proporcional};
\node[align=center,font=\scriptsize] at (7.25,-1.2){prop.\\limitada};
\node[align=center,font=\scriptsize] at (8.95,-1.2){Geiger--\\Müller};
\node[align=center,font=\scriptsize] at (0.7,-1.2){recomb.};
\end{tikzpicture}
\end{center}
\begin{enumerate}[label=(\alph*)]
\item \textbf{Recombinación}: campo insuficiente; parte de los pares se recombina.
\item \textbf{Cámara de ionización}: se recogen todos los pares primarios; meseta; señal $\propto$ energía e independiente de $V$.
\item \textbf{Contador proporcional}: los electrones ganan energía suficiente para ionizar (avalancha de Townsend); carga
$= M\times$ primarios, $M$ crece exponencialmente con $V$. Sigue siendo $\propto$ energía.
\item \textbf{Proporcionalidad limitada}: la carga espacial de los iones distorsiona el campo; se pierde la proporcionalidad.
\item \textbf{Geiger--Müller}: descarga en todo el hilo; misma señal para cualquier partícula: sólo cuenta. Después, descarga continua.
\end{enumerate}
El tipo de detector lo determina principalmente la tensión aplicada.
\end{clave}

\begin{clave}[Tamaño de la señal: por qué hace falta multiplicación]
\begin{itemize}
\item $\alpha$ de 5.5 MeV parada en Ar: $N=5.5\cdot10^6/26\approx2.1\cdot10^5$ pares $\Rightarrow Q\approx34$ fC $\Rightarrow\Delta V=Q/C\approx3.4$ mV sobre 10 pF.
\item $\beta$ rápido que cruza unos cm de gas: $\sim10$ keV $\to\sim400$ pares $\to Q\approx0.06$ fC: demasiado pequeño para contar
individualmente $\Rightarrow$ modo corriente o multiplicación.
\item Proporcional con $M=10^3$: $Q\approx60$ fC, fácilmente medible.
\item GM: $10^9$--$10^{10}$ electrones por descarga $\to$ pulsos de voltios, sin información de energía.
\end{itemize}
\end{clave}

\section{Cámaras de ionización}
\begin{itemize}
\item Gas: He, N$_2$, Ar, Kr, Xe\dots Geometrías variadas (placas paralelas, cilindros).
\item Amplitud $\propto$ número de iones $\propto$ energía depositada, independiente de $V$.
\item Eficiencia $\sim100\%$ para cargadas, $\sim1\%$ para fotones (baja densidad del gas).
\item Señales muy pequeñas ($\sim10^{-4}$ V) y lentas (la deriva de iones es lenta, hasta $\sim$ms) $\Rightarrow$ gran amplificación.
\item Usos: dosimetría y monitores de radiación (modo corriente), detectores de humo (fuente $\alpha$ de $^{241}$Am; el humo
reduce la corriente), etapa $\Delta E$ en telescopios $\Delta E$(gas)--$E$(Si) para iones pesados.
\end{itemize}

\section{Contadores proporcionales}
\begin{itemize}
\item Ánodo de hilo fino en un cátodo cilíndrico. El campo es
\[ E(r)=\frac{V}{r\ln(b/a)} \]
($a$ radio del hilo, $b$ del cátodo): sólo muy cerca del hilo es suficientemente intenso para la avalancha.
Así la multiplicación es independiente del punto de interacción y la señal sigue siendo proporcional.
\item Factor de multiplicación (ganancia gaseosa) $M\sim10^2$--$10^4$, dependencia \textbf{exponencial} con $V$
$\Rightarrow$ fuente de alta tensión muy estable.
\item Gases: Ar + gas de extinción (P10 = 90\% Ar + 10\% CH$_4$), isobutano, $^3$He o BF$_3$ (neutrones, Tema 6).
\item Modo pulso, respuesta $\sim\mu$s. Permite espectroscopía de baja energía (X blandos, $\beta$).
\end{itemize}

\subsection{Detectores sensibles a posición}
\begin{itemize}
\item \textbf{MWPC} (cámara proporcional multihilos, Charpak 1968, Nobel 1992): plano de hilos anódicos finos
($\sim0.1$ mm de diámetro... separados $\sim$mm) entre dos planos catódicos. Cada hilo actúa como un contador
proporcional independiente: da energía y posición ($X,Y$).
\item \textbf{Cámaras de deriva y TPC}: la posición se obtiene del tiempo de deriva de los electrones.
\item \textbf{MPGD} (micropatrón): \textbf{GEM} (lámina de Kapton perforada, $\sim$400 V, ganancia 10--100 por lámina, en cascadas de 3)
y \textbf{Micromegas} (micromalla a 50--130 $\mu$m del ánodo). Resolución $\sim$50--100 $\mu$m, tasas $>10^6$ Hz/mm$^2$, bajo retorno de iones.
Uso: TPC de ALICE, TPC de blanco activo en física nuclear.
\end{itemize}
\begin{dato}
Con MWPC se descubrieron el quark charm (Richter y Ting, 1974; Nobel 1976) y los bosones W y Z (Rubbia y van der Meer, 1983; Nobel 1984).
\end{dato}

\section{Contadores Geiger--Müller}
\begin{itemize}
\item Gas He, Ne o Ar (+ gas de extinción o \emph{quenching}, orgánico o halógeno); HV 400--1200 V.
\item La avalancha se propaga por todo el hilo (fotones UV de desexcitación): pulso grande e igual para todas las partículas.
\item \textbf{No permite espectroscopía}: sólo contar ($\beta$, $\gamma$). Ventana de vidrio (robusta, no ve $\beta$ blandos) o de mica (frágil, permite $\beta$/$\alpha$).
\item Tiempo muerto grande (decenas--cientos de $\mu$s): la nube de iones positivos tarda en llegar al cátodo; no paralizable en la práctica.
\end{itemize}

\begin{center}\small
\begin{tabularx}{\linewidth}{lXXX}
\toprule
& Cámara de ionización & Proporcional & Geiger--Müller\\
\midrule
Multiplicación & No ($M=1$) & Sí, $M\sim10^2$--$10^4$ & Descarga total\\
Señal $\propto E$ & Sí & Sí & No\\
Tamaño de señal & Muy pequeño & Medio & Grande (V)\\
Uso típico & Dosimetría, $\Delta E$ & Espectroscopía X, neutrones, trazas & Contaje, protección radiológica\\
\bottomrule
\end{tabularx}
\end{center}

\chapter{Detectores de semiconductor}

\section{Idea básica}
Un detector de semiconductor es una ``cámara de ionización de estado sólido'': volumen con campo
eléctrico donde la energía depositada crea \textbf{pares electrón--hueco} que derivan hacia los
electrodos e inducen una corriente en el circuito externo.

\begin{clave}[Ventajas frente a gas]
\begin{itemize}
\item \textbf{Energía por par muy baja}: 3.6 eV en Si, 2.96 eV en Ge (vs $\approx30$ eV en gas) $\Rightarrow$ muchos
portadores $\Rightarrow$ \textbf{excelente resolución en energía}.
\item \textbf{Densidad alta} (sólido): menor alcance para la misma energía, detectores compactos.
\item Campos altos y alta movilidad: \textbf{respuesta rápida} y buena eficiencia de colección (poco atrapamiento).
\item Resolución en posición dada por el tamaño de los electrodos (segmentación).
\end{itemize}
\end{clave}
\begin{dato}[gas / líquido / sólido]
Densidad: baja / moderada / alta. $Z$: bajo / moderado / moderado. Energía de ionización: moderada / moderada / baja.
Velocidad de señal: moderada / moderada / rápida.
\end{dato}

\section{Física del semiconductor}
\begin{itemize}
\item Si y Ge: 4 electrones de valencia, red tipo diamante. Gap: Si 1.12 eV, Ge 0.67 eV ($k_BT\approx25$--40 meV a temperatura ambiente).
\item \textbf{Dopado tipo n}: impureza pentavalente (P, As): donores, electrones mayoritarios.
\textbf{Tipo p}: impureza trivalente (B): aceptores, huecos mayoritarios.
\item Problema: en Si intrínseco $n_i\approx1.45\cdot10^{10}$ cm$^{-3}$ $\Rightarrow$ en un volumen típico hay $\sim4.5\cdot10^8$ portadores libres frente
a sólo $3.2\cdot10^4$ pares creados por una MIP. \textbf{Hay que eliminar los portadores libres: deplexionar} el detector
con una unión p--n polarizada en inversa.
\end{itemize}

\subsection{Unión p--n}
Por difusión, electrones del lado n pasan al p y huecos al revés, dejando carga espacial fija
(donores ionizados $+$ en n, aceptores $-$ en p). Se forma una \textbf{zona de deplexión} sin portadores
libres y una barrera de potencial que frena la difusión.
\begin{itemize}
\item \textbf{Polarización directa} ($+$ a p): baja la barrera, gran corriente.
\item \textbf{Polarización inversa} ($+$ a n): aumenta la barrera y la zona deplexionada; sólo circula una pequeña
\textbf{corriente de fuga} (portadores minoritarios), que crece con $T$, el espesor y la irradiación: es ruido.
\end{itemize}

\begin{clave}[Zona de deplexión y tensión de deplexión total]
Ecuación de Poisson (unión abrupta p$^+$--n):
\[ -\frac{\mathrm{d}^2\phi}{\mathrm{d}x^2}=\frac{q_0}{\varepsilon\varepsilon_0}N_{eff} \]
La densidad de carga es constante en la zona deplexionada, el campo crece linealmente (máximo en la unión) y el potencial es parabólico.
\[
w=\sqrt{\frac{2\varepsilon\varepsilon_0}{q_0|N_{eff}|}V}\ \propto\sqrt{V_B},\qquad
V_{dep}=\frac{q_0}{\varepsilon\varepsilon_0}\,|N_{eff}|\,d^2\quad\left(\text{exacto: }\tfrac{q_0}{2\varepsilon\varepsilon_0}|N_{eff}|d^2\right)
\]
\begin{itemize}
\item La zona deplexionada \textbf{crece con la tensión} ($w\propto\sqrt{V_B}$).
\item Los pares creados en la zona \emph{no} deplexionada (sin campo) no dan señal:
\textbf{colección completa de carga sólo si el detector está totalmente deplexionado ($V_B>V_{dep}$)}.
\item $V_{dep}\propto N_{eff}\,d^2$: material más puro (alta resistividad) o más fino $\Rightarrow$ menor tensión.
\end{itemize}
\end{clave}
\begin{dato}
Si tipo n de $\rho=2$ k$\Omega$cm ($N_D\approx2.2\cdot10^{12}$ cm$^{-3}$), 300 $\mu$m: $V_{dep}\approx150$ V, fuga $<1\,\mu$A.
\end{dato}

\section{Detectores de silicio}
\begin{clave}[Si: principalmente partículas cargadas y rayos X]
\begin{itemize}
\item Gap 1.12 eV, 3.6 eV/par: funciona a temperatura ambiente.
\item $\rho=2.33$ g/cm$^3$; MIP: $\dEdx\approx3.8$ MeV/cm $\approx108$ pares/$\mu$m.
\item Alta movilidad ($\mu_e=1450$, $\mu_h=450$ cm$^2$/Vs) $\Rightarrow$ colección $<10$ ns. Velocidad de deriva $v=\mu E$.
\item Muy puro, rígido (estructuras finas autoportantes), tecnología microelectrónica barata.
\item $Z=14$ bajo: poco eficiente para $\gamma$ de alta energía.
\end{itemize}
\end{clave}

\subsection{Silicios segmentados}
Segmentando el contacto p$^+$ en \textbf{tiras} (\emph{strips}) leídas por canales independientes se obtiene la posición
(resolución $\propto$ paso de las tiras, $\sigma\approx\text{paso}/\sqrt{12}$). Con tiras perpendiculares en ambas caras
(\textbf{DSSD}, \emph{double-sided strip detector}) se obtiene $(x,y)$. Espesor típico \textbf{300 $\mu$m} (150--500 $\mu$m).
Geometrías: cuadrados, anulares (tipo CD), en cajas alrededor del blanco; se combinan con Ge para medir partículas y $\gamma$ en coincidencia.

\subsection{Telescopios $\Delta E$--$E$ e identificación de partículas}
\begin{clave}[Telescopio $\Delta E$--$E$]
Un detector \textbf{muy delgado (5--50 $\mu$m)} mide la pérdida $\Delta E$ y otro \textbf{grueso (200--1500 $\mu$m)} detrás
para completamente la partícula ($E$ residual). \textbf{Sumando ambas señales se obtiene la energía total}.

Por \textbf{Bethe--Bloch}, en régimen no relativista
\[ \Delta E\propto\frac{Mz^2}{E}\quad\Rightarrow\quad \Delta E\cdot E\propto Mz^2 \]
cada especie (p, d, t, $^3$He, $\alpha$\dots) ocupa una \textbf{hipérbola distinta} en el diagrama $\Delta E$--$E$:
identificación de partículas. Seleccionando una banda se aísla un canal de reacción (p.ej.\ $^{163}$Dy($^3$He,$^4$He)$^{162}$Dy).
\end{clave}

\subsection{Relación señal/ruido}
La distribución de pérdida en capas finas (Landau) tiene cola a baja energía, ensanchada por el ruido. Fuentes de ruido
(ENC, carga equivalente de ruido): capacidad, corriente de fuga, ruido térmico de la resistencia de polarización.
Se exige un umbral por encima de la cola de ruido (alto $\Rightarrow$ pérdida de eficiencia; bajo $\Rightarrow$ ocupación por ruido).
Figura de mérito: $S/N>10$--15. El daño por radiación la degrada.

\section{Detectores de germanio (HPGe)}
\begin{clave}[HPGe: espectroscopía $\gamma$ y X de alta resolución]
\begin{itemize}
\item $Z=32$ (vs 14 del Si) y cristales grandes (coaxiales, cientos de cm$^3$): adecuado para $\gamma$.
\item Gap pequeño (0.67 eV) $\Rightarrow$ \textbf{hay que enfriar a 77 K (N$_2$ líquido)} o con refrigeración eléctrica para suprimir la
generación térmica de pares (corriente de fuga).
\item Alta pureza ($\sim10^{9}$--$10^{10}$ impurezas/cm$^3$, frente a $10^{12}$ en los antiguos Ge(Li), que necesitaban estar siempre fríos):
permite deplexionar varios cm con $V>1000$ V.
\item \textbf{Resolución $\approx0.15\%$ a 1.33 MeV} (FWHM 1.7--2.3 keV) frente a $\approx7$--8\% del NaI. Mejor cociente pico/Compton.
\item \textbf{Eficiencia intrínseca menor} que la de los centelleadores inorgánicos de $Z$ y $\rho$ altos.
\end{itemize}
\end{clave}
\begin{itemize}
\item Geometrías planar y coaxial; contactos de Li difundido (n$^+$, $\sim0.5$ mm, capa muerta gruesa) o B implantado (p$^+$, $\sim0.5\,\mu$m).
\item Para X y $\gamma$ de baja energía la ventana de entrada debe ser fina: contacto p$^+$ delante $\Rightarrow$ cristal tipo n
(además, el tipo n es más resistente al daño por neutrones). Ventana de Be para bajas energías, de Al para altas.
\item Cadena: detector $\to$ preamplificador $\to$ amplificador $\to$ ADC $\to$ ordenador. Constantes de conformado 2--6 $\mu$s.
\end{itemize}
\begin{dato}[Parámetros estándar a 1332 keV ($^{60}$Co)]
Eficiencia relativa (respecto a NaI 3''$\times$3'' a 25 cm, $\varepsilon=1.2\cdot10^{-3}$): 10--70\% (hoy $>100\%$).
Pico/Compton 30--60. FWHM 1.7--2.3 keV. Forma: FWTM/FWHM $\approx2.0$ (gaussiana 1.82).
\end{dato}

\begin{clave}[Detectores HPGe compuestos y segmentados]
\begin{itemize}
\item Para \textbf{mejorar el cociente pico/Compton} y \textbf{corregir el desplazamiento Doppler} (fuentes en vuelo) hay que conocer
dónde interacciona el $\gamma$: \textbf{tracking} $\Rightarrow$ cristales de Ge \textbf{segmentados} + análisis de forma de pulso.
\item Para aumentar la eficiencia se agrupan en \textbf{arrays} (clústeres con criostato común; detectores \emph{clover}).
\item Supresión Compton: escudos de BGO alrededor que vetan los sucesos en que el fotón escapa.
\end{itemize}
\end{clave}
\begin{dato}[Ejemplos]
Euroball, Exogam (clovers), Miniball (ISOLDE, cristales segmentados en 6), \textbf{AGATA} y \textbf{GRETINA/GRETA}
(tracking: 36 segmentos por cristal, 6 longitudinales $\times$ 6 transversales).
\end{dato}

\section{Otros semiconductores}
CdTe, CdZnTe (CZT), HgI$_2$: $Z$ alto y gap mayor $\Rightarrow$ \textbf{no requieren refrigeración} ($\varepsilon\approx4.4$ eV/par),
pero peor transporte de carga (huecos); por ahora para energías bajas. Diamante, GaAs, SiC: resistentes a la radiación.
SiPM (fotomultiplicadores de silicio): ver Tema 5.

\chapter{Detectores de centelleo}

\section{Principio}
\begin{clave}[Detección indirecta]
La radiación excita átomos/moléculas del centelleador; al desexcitarse emiten luz (UV/visible), que se
guía hasta un \textbf{fotosensor} (PMT, APD, SiPM) que la convierte en señal eléctrica y la amplifica.
\[ \text{centelleador}\ \to\ \text{guía de luz / WLS}\ \to\ \text{fotosensor}\ \to\ \text{pulso eléctrico} \]
La cantidad de luz crece con la energía depositada $\Rightarrow$ pueden funcionar como espectrómetros en modo pulso.
\end{clave}

\paragraph{Propiedades deseables:} alta eficiencia de conversión (rendimiento luminoso, mejor resolución), tiempos de caída cortos,
transparencia a su propia emisión, espectro de emisión adaptado al fotosensor, $\rho$ y $Z$ altos para $\gamma$, robustez,
no higroscópicos, baratos. Emisiones: fluorescencia (inmediata), fluorescencia retardada y fosforescencia (retardada, distinta $\lambda$).

Aplicaciones médicas: gammacámaras, TC, PET (BGO; PET-TOF moderno con LSO/LYSO), radiografía digital.

\section{Centelleadores inorgánicos y orgánicos}
\begin{clave}[Inorgánicos vs orgánicos]
\begin{center}\small
\begin{tabularx}{\linewidth}{lXX}
\toprule
& \textbf{Inorgánicos} (cristales) & \textbf{Orgánicos} (cristales, líquidos, plásticos)\\
\midrule
Mecanismo & Estructura de bandas de la red cristalina + \textbf{activadores} (impurezas) & Transiciones \textbf{moleculares} ($\pi$-electrones), independiente del estado físico\\
$Z$, $\rho$ & Altos $\Rightarrow$ \textbf{radiación $\gamma$} & Bajos (C, H) $\Rightarrow$ domina Compton; buenos para cargadas, $\beta$ y \textbf{neutrones rápidos} (H)\\
Luz & Mucha (NaI 38\,000 fot/MeV) & Menos ($\sim$10\,000 fot/MeV)\\
Velocidad & Lentos (decenas--cientos de ns) & \textbf{Muy rápidos} ($\sim$2 ns)\\
Ejemplos & NaI(Tl), CsI(Tl), BGO, BaF$_2$, LaBr$_3$(Ce), LSO, LYSO, PbWO$_4$ & Antraceno, estilbeno, líquidos (NE213), plásticos (BC-400, NE111)\\
\bottomrule
\end{tabularx}
\end{center}
\end{clave}

\subsection{Luminiscencia en cristales inorgánicos}
En un cristal puro, el electrón excitado a la banda de conducción vuelve a la de valencia emitiendo un fotón de energía del gap:
poco probable y en el UV, y el cristal lo reabsorbe. Se añaden \textbf{activadores} (Tl, Ce\dots) que crean niveles dentro del gap:
el par $e$--$h$ migra hasta el activador, que se desexcita emitiendo un fotón de \textbf{menor energía} (visible) que
\textbf{no es reabsorbido} por el cristal. Otros mecanismos: excitones, defectos, luminiscencia autoactivada (BGO) y
luminiscencia cruzada (componente rápida del BaF$_2$).

\subsection{Luminiscencia en orgánicos}
La radiación excita los estados singlete $S_1$ (y vibracionales superiores), que decaen sin radiar al fondo de $S_1$ y de ahí a niveles
vibracionales de $S_0$ emitiendo luz (fluorescencia rápida, ns). Como la emisión parte del fondo de $S_1$ y acaba en niveles
vibracionales excitados de $S_0$, el espectro de emisión está desplazado a $\lambda$ mayores que el de absorción
$\Rightarrow$ \textbf{transparente a su propia luz}. Los estados triplete (cruce intersistema) dan componentes lentas (fosforescencia/fluorescencia retardada).

\paragraph{Forma del pulso.}
\[ I(t)=I_0\left(e^{-t/\tau_f}-e^{-t/\tau_r}\right) \]
$\tau_r\approx0.2$--0.4 ns (transferencia al estado fluorescente), $\tau_f\approx1$--3 ns (decaimiento) en orgánicos.

\subsection{Ley de Birks (quenching de ionización)}
\begin{clave}[Ley de Birks]
\[
\frac{\mathrm{d}L}{\mathrm{d}r}=\frac{S\,\dfrac{\mathrm{d}E}{\mathrm{d}r}}{1+kB\,\dfrac{\mathrm{d}E}{\mathrm{d}r}}
\]
\begin{itemize}
\item $\mathrm{d}E/\mathrm{d}r$ pequeño (electrones): $\mathrm{d}L/\mathrm{d}r\approx S\,\mathrm{d}E/\mathrm{d}r$, luz proporcional a la energía.
\item $\mathrm{d}E/\mathrm{d}r$ grande (p, $\alpha$, iones): \textbf{saturación}, $\mathrm{d}L/\mathrm{d}r\to S/kB$: menos luz por MeV.
\item Consecuencia: la luz producida y la \textbf{forma del pulso} dependen del tipo de partícula $\Rightarrow$
\textbf{discriminación por forma de pulso (PSD)}, p.ej.\ neutrones/$\gamma$ en centelleadores orgánicos líquidos; CsI(Tl) para p/$\alpha$.
Por eso la energía se suele dar en keV equivalentes de electrón (keVee).
\end{itemize}
\end{clave}

\section{Principales centelleadores}
\begin{clave}[Tabla de centelleadores (del resumen oficial)]
\begin{center}\small
\begin{tabular}{lcccccc}
\toprule
& $\rho$ (g/cm$^3$) & $Z_{eff}$ & $\lambda_{max}$ (nm) & $n$ & $\tau$ (ns) & Luz (fot/MeV)\\
\midrule
NaI(Tl) & 3.67 & 51 & 415 & 1.85 & 230 & 38\,000\\
CsI(Tl) & 4.51 & 54 & 540 & 1.80 & 680, 3340 & 40\,000, 25\,000\\
BGO (Bi$_4$Ge$_3$O$_{12}$) & 7.13 & 75 & 480 & 2.15 & 300 & 8\,200\\
BaF$_2$ & 4.89 & 53 & 220, 310 & 1.56 & 0.6, 630 & 1\,500, 9\,500\\
CeF$_3$ & 6.16 & 53 & 310, 340 & 1.68 & 5, 27 & 4\,400\\
YAlO$_3$(Ce) & 5.37 & 34 & 370 & 1.95 & 27 & 18\,000\\
LSO (Lu$_2$SiO$_5$:Ce) & 7.4 & 66 & 420 & 1.82 & 47 & 25\,000\\
LaBr$_3$(Ce) & 3.79 & 47 & 350 & 1.9 & 25 & 63\,000\\
GSO & 6.71 & 59 & 440 & 1.85 & 60 & 7\,500\\
LYSO & 7.1 & 65 & 420 & 1.81 & 40 & 41\,000\\
Plástico BC-400 & 1.03 & 4 & 420 & 1.58 & 2.4 & 10\,000\\
\bottomrule
\end{tabular}
\end{center}
\textbf{Qué retener} (más que los números exactos):
\begin{itemize}
\item \textbf{NaI(Tl)}: el estándar; mucha luz, barato, grande; \emph{higroscópico y frágil}; 230 ns. Medicina nuclear.
\item \textbf{BGO}: $\rho$ y $Z$ máximos $\Rightarrow$ máxima eficiencia, pero poca luz (mala resolución). PET clásico, escudos anti-Compton.
\item \textbf{BaF$_2$}: componente ultrarrápida (0.6 ns, UV) $\Rightarrow$ timing.
\item \textbf{LaBr$_3$(Ce)}: máximo rendimiento luminoso y rápido $\Rightarrow$ mejor resolución entre centelleadores ($\approx$3\% a 662 keV) y timing (vidas medias, Tema 11).
\item \textbf{LSO/LYSO}: densos y rápidos $\Rightarrow$ PET-TOF.
\item \textbf{Plásticos}: muy rápidos, baratos, grandes; ricos en H (neutrones, ToF); $Z$ bajo.
\end{itemize}
\end{clave}
\begin{dato}
Plásticos: NE111 con subida 0.2 ns y caída 1.7 ns. Los orgánicos líquidos son resistentes al daño por radiación y la muestra
radiactiva puede disolverse en ellos (centelleo líquido). Añadir Pb a un plástico aumenta el fotopico pero reduce la luz.
PbWO$_4$ (CMS, ALICE): muy denso y rápido pero poca luz. La luz de muchos inorgánicos depende de la temperatura.
\end{dato}

\section{Colección de la luz}
Ley de Snell $n_i\sin\theta_i=n_t\sin\theta_t$ (reflexión total interna) y fórmulas de Fresnel. Se envuelve el cristal con
reflector (especular o difuso, p.ej.\ teflón), se acopla con grasa óptica (contacto óptico) y, si hace falta, guías de luz o
fibras desplazadoras de longitud de onda (WLS).

\section{Fotosensores}
\subsection{Tubo fotomultiplicador (PMT)}
\begin{clave}[PMT]
Tubo de vidrio al vacío con \textbf{fotocátodo} (efecto fotoeléctrico, eficiencia cuántica QE $\approx25\%$), electrodo de
focalización, 10--12 \textbf{dínodos} y ánodo. Cada electrón que llega a un dínodo arranca $\delta\approx5$ secundarios:
\[ G=\delta^{\,n}\qquad(\delta=5,\ n=10\ \Rightarrow\ G\approx10^7) \]
\begin{itemize}
\item Ganancia $10^6$--$10^7$, se ajusta con la tensión (\textbf{$\sim$1 kV}); muy dependiente de la HV.
\item Gran área, muy bajo ruido, rápido (200--500 ps).
\item \textbf{Sensible a campos magnéticos} (desvían los electrones entre dínodos).
\end{itemize}
\end{clave}

\subsection{Resolución de un centelleador}
\begin{clave}[El cuello de botella son los fotoelectrones]
\[
N_{pe}=E\cdot LY\cdot\eta_{col}\cdot QE,\qquad R\approx2.355\sqrt{\frac{1+\nu}{N_{pe}}}
\]
($\nu\approx0.1$--0.2: varianza relativa de la ganancia del PMT).
NaI(Tl) a 662 keV: $38\,000\times0.662\approx25\,000$ fotones; $\eta_{col}\approx0.5$, $QE\approx0.25$
$\Rightarrow N_{pe}\approx3100\Rightarrow R\approx4.5\%$ (estadístico). Medido $\approx6.5$--7\% (no proporcionalidad intrínseca,
inhomogeneidades). LaBr$_3$: $\approx2.6$--3\%. HPGe: $\approx2.2\cdot10^5$ pares $\Rightarrow\approx0.2\%$.
\end{clave}

\subsection{Fotosensores de estado sólido y otros}
\begin{itemize}
\item \textbf{Fotodiodo de avalancha (APD)}: unión polarizada con tensión alta (100--500 V, hasta 2 kV): el fotoelectrón crea más pares
(modo proporcional). Ganancia moderada y \emph{muy dependiente de la temperatura}; mayor QE; \textbf{insensible a campos magnéticos}; compacto.
\item \textbf{SiPM} (G-APD, MPPC): \textbf{matriz} de micro-APD (20--100 $\mu$m, $\sim$1000/mm$^2$) en \textbf{modo Geiger} (por encima de la tensión
de ruptura). Cada celda da una señal fija; la salida es la suma $\Rightarrow$ cuenta fotones. Ganancia $\sim10^6$, \textbf{baja tensión (20--60 V)},
rápidos, compactos, baratos, \textbf{inmunes a campos magnéticos}. Inconvenientes: ruido oscuro alto (térmico), saturación, PDE 20--50\%.
\item \textbf{MCP} (placas de microcanales): capilares emisores de electrones secundarios; resolución temporal (20--100 ps) y espacial (10--50 $\mu$m) excelentes;
área pequeña, vida limitada, caros.
\end{itemize}
\begin{clave}[Del resumen]
\textbf{PMTs sensibles a campos magnéticos y SiPMs no. PMTs requieren voltajes del orden de 1 kV} (SiPM: decenas de V).
\end{clave}

\section{Centelleador NaI frente a semiconductor HPGe}
\begin{clave}[Comparación (espectros de $^{137}$Cs y $^{60}$Co)]
\begin{itemize}
\item \textbf{HPGe: mejor resolución en energía} (FWHM $\sim$2 keV vs $\sim$50 keV a 662 keV): picos muy estrechos, separa líneas próximas.
\item \textbf{NaI: mejor eficiencia intrínseca} ($Z$ del I, cristales grandes y baratos).
\item \textbf{Los centelleadores son en general más rápidos}.
\item NaI funciona a temperatura ambiente; HPGe necesita refrigeración.
\end{itemize}
\end{clave}

\chapter{Detección de neutrones}

\section{Idea básica}
\begin{clave}[Detección siempre indirecta]
Los neutrones no tienen carga: no ionizan. Hay que usar \textbf{reacciones nucleares} que los ``conviertan'' en partículas
cargadas (o $\gamma$), que luego se detectan con detectores convencionales (gas, centelleador, semiconductor).
\begin{itemize}
\item \textbf{Neutrones lentos/térmicos}: reacciones de \textbf{captura} exotérmicas con gran sección eficaz.
\item \textbf{Neutrones rápidos}: \textbf{dispersión elástica} (protones de retroceso en materiales hidrogenados) o se moderan primero.
\end{itemize}
\end{clave}

\paragraph{Canales de interacción:} dispersión elástica (n,n), inelástica (n,n'$\gamma$), captura radiativa (n,$\gamma$), multiplicación
(n,xn), producción de cargadas (n,p), (n,$\alpha$), y fisión (n,f). Sin barrera coulombiana; algunas con umbral.
Dependencia con la energía: región \textbf{$1/v$} a baja energía, \textbf{resonancias} y continuo a alta energía.

\begin{clave}[Reacciones de conversión (secciones eficaces térmicas, 0.025 eV)]
\begin{center}
\begin{tabular}{lcc}
\toprule
Reacción & $Q$ (MeV) & $\sigma_{th}$ (b)\\
\midrule
n + $^3$He $\to$ $^3$H + p & 0.764 & 5330\\
n + $^{10}$B $\to$ $^7$Li$^*$ + $\alpha$ $\to$ $^7$Li + $\alpha$ + $\gamma$(0.48) & 2.31 (93\%) / 2.79 (7\%) & 3840\\
n + $^6$Li $\to$ $^3$H + $\alpha$ & 4.78 & 940\\
n + $^{155,157}$Gd $\to$ Gd$^*$ $\to\gamma$ + e$^-$ de conversión & -- & $\sim10^5$\\
n + $^{235}$U, $^{239}$Pu $\to$ fisión & $\sim$160--200 & 585 / 750\\
\bottomrule
\end{tabular}
\end{center}
Todas siguen la ley \textbf{$\sigma\propto1/v$} a baja energía $\Rightarrow$ eficientes sólo para neutrones lentos: para neutrones rápidos hay
que \textbf{moderar} primero. $Q$ grande facilita discriminar frente a $\gamma$ (pulsos más grandes que el fondo).
\end{clave}

\section{Clasificación}
\begin{itemize}
\item \textbf{Contadores} (sólo detectan): moderados o no moderados.
\item \textbf{Espectrómetros} (energía): retroceso, reacción con partícula cargada, \textbf{tiempo de vuelo}, moderación (esferas de Bonner).
\item Por forma física: gas (cámaras de ionización y proporcionales), líquido (centelleadores), sólido (centelleadores, semiconductores).
\item Por material activo: autodetector, cargado (con $^{10}$B o $^6$Li añadido), recubierto (capa convertidora).
\end{itemize}

\section{Detectores gaseosos}
\begin{clave}[Contadores proporcionales de $^3$He y BF$_3$]
\begin{itemize}
\item $^3$He(n,p)$^3$H: 0.764 MeV repartidos entre p (573 keV) y t (191 keV). $\sim25\,000$ pares por neutrón ($\sim4\cdot10^{-15}$ C).
En modo proporcional, ganancias de hasta miles.
\item Espectro térmico: pico a 764 keV + \textbf{efecto pared} (una de las partículas toca la pared antes de depositar toda su energía):
continuo con escalones en 191 y 573 keV.
\item Muy buena discriminación $n/\gamma$ por altura de pulso (los $\gamma$ depositan poco en el gas).
\item BF$_3$: gas tóxico, menor $\sigma$; $^{10}$B(n,$\alpha$).
\item Otros gases: H$_2$, CH$_4$, $^4$He (retroceso, neutrones rápidos).
\end{itemize}
\end{clave}
\begin{dato}[Crisis del $^3$He]
El $^3$He procede de la desintegración del tritio (arsenal nuclear): escasez desde $\sim$2009. Alternativas: BF$_3$, tubos recubiertos de $^{10}$B,
láminas de $^{10}$B$_4$C (Multi-Grid, ESS), vidrios de $^6$Li, $^6$LiF:ZnS(Ag), elpasolitas (CLYC) con PSD.
\end{dato}

\section{Centelleadores}
\begin{itemize}
\item \textbf{Con $^6$Li} (neutrones lentos): vidrio Li(Ce), LiI(Eu), ZnS(Ag)--LiF (mucha luz por neutrón, opaco).
\item \textbf{Orgánicos líquidos} (BC501/NE213) y plásticos (neutrones rápidos): protones de retroceso; espectro casi rectangular entre 0 y $E_n$
(distorsionado por la no linealidad de Birks y la resolución). Permiten \textbf{PSD} n/$\gamma$ (la componente lenta es mayor para protones).
\item CLYC (Cs$_2$LiYCl$_6$:Ce): neutrones térmicos ($^6$Li) y rápidos ($^{35}$Cl(n,p)) con PSD.
\end{itemize}

\section{Semiconductores}
$\sim1.5\cdot10^6$ pares por neutrón: se detecta sin amplificación adicional. Pero los semiconductores convencionales no tienen
núcleos absorbentes $\Rightarrow$ se deposita una \textbf{capa convertidora} ($^6$LiF, $^{10}$B) sobre el Si. La capa debe ser fina (pocos $\mu$m)
para que salgan las partículas cargadas $\Rightarrow$ \textbf{eficiencia baja} (y sólo una de las dos partículas entra en el detector).

\section{Dispersión elástica y moderación}
\begin{clave}[Moderación]
Tras una colisión elástica (isótropa en el CM, onda $s$) la energía del neutrón queda uniformemente entre:
\[
\alpha E_0\le E\le E_0,\qquad \alpha=\left(\frac{A-1}{A+1}\right)^2
\]
\begin{itemize}
\item \textbf{Buen moderador: material ligero rico en hidrógeno} ($\alpha=0$: el neutrón puede perder toda su energía en un choque).
Polietileno, agua, parafina.
\item \textbf{Letargia}: $u=\ln(E_0/E)$. Decremento logarítmico medio por colisión:
\[ \xi=\left\langle\ln\frac{E_0}{E}\right\rangle=1+\frac{\alpha}{1-\alpha}\ln\alpha=1-\frac{(A-1)^2}{2A}\ln\frac{A+1}{A-1} \]
\item Número de colisiones para pasar de $E_0$ a $E_f$:
\[ N=\frac{\ln(E_0/E_f)}{\xi} \]
\end{itemize}
\begin{center}\small
\begin{tabular}{lccc}
\toprule
Núcleo & $1-\alpha$ & $\xi$ & $N$ (1 MeV $\to$ 25 meV)\\
\midrule
$^1$H & 1 & 1 & 18\\ $^2$H & 0.889 & 0.725 & 24\\ $^4$He & 0.640 & 0.425 & 41\\
$^{12}$C & 0.284 & 0.158 & 111\\ $^{56}$Fe & 0.069 & 0.035 & 500\\ $^{208}$Pb & 0.019 & 0.010 & 1823\\
\bottomrule
\end{tabular}
\end{center}
Al moderar, el neutrón pasa a la región $1/v$ donde las secciones eficaces de captura ($^3$He, $^{10}$B, $^6$Li) son enormes.
\end{clave}
\begin{ejemplo}
De 2 MeV a 0.025 eV en hidrógeno: $N=\ln(8\cdot10^7)/1\approx18$ colisiones.
\end{ejemplo}

\section{Espectrómetro de tiempo de vuelo (ToF)}
\begin{clave}[ToF]
\[ E=\frac12mv^2=\frac12m\frac{L^2}{t^2},\qquad \frac{\Delta E}{E}=2\sqrt{\left(\frac{\Delta L}{L}\right)^2+\left(\frac{\Delta t}{t}\right)^2} \]
\begin{itemize}
\item \textbf{Start}: origen pulsado en el tiempo (haz pulsado) o radiación acompañante ($\beta$, $\alpha$, $\gamma$\dots, \emph{no} el neutrón).
\item \textbf{Stop}: el detector de neutrones.
\item \textbf{Base de vuelo larga, detectores cortos (poca incertidumbre $\Delta L$ en el punto de interacción), buena resolución temporal}.
\item Detectores típicos: centelleadores plásticos/orgánicos (rápidos, ricos en H).
\end{itemize}
\end{clave}
\begin{ejemplo}[$L=1$ m]
$E=1$ MeV: $\beta=0.046$, $t\approx72$ ns (regla: $t[\mathrm{ns}]\approx72.3\,L[\mathrm{m}]/\sqrt{E[\mathrm{MeV}]}$).
Con $\Delta t=1$ ns y $\Delta L=1$ cm: $\Delta E/E=2\sqrt{0.01^2+(1/72)^2}\approx3.4\%$. A 10 MeV: $t\approx23$ ns, $\Delta E/E\approx9\%$.
Neutrones rápidos $\Rightarrow$ bases largas (n\_TOF en el CERN: 185 m y $\sim$20 m).
\end{ejemplo}
\begin{dato}[Instrumentos]
\textbf{TONERRE} (LPC-Caen, ISOLDE): centelleador plástico BC400, $\Delta\Omega/4\pi=50\%$, $\varepsilon_{int}=25\%$, $\Delta E/E=10\%$ (neutrones $\beta$-retardados).
\textbf{MoNA} (NSCL): 144 barras plásticas de $200\times10\times10$ cm$^3$. \textbf{Neutron Wall} (EUROBALL): líquido BC501, $\Delta\Omega/4\pi=25\%$, $\varepsilon_{int}=50\%$.
\end{dato}

\section{Detectores moderados}
\begin{clave}[Esferas de Bonner]
Esferas de polietileno de distintos diámetros (2''--18'') con un contador de neutrones térmicos en el centro ($^3$He o $^6$LiI).
Cada esfera tiene una \textbf{función respuesta} $R_d(E_n)$ distinta (calculada por Monte Carlo): las pequeñas responden a neutrones
lentos, las grandes a rápidos (se añaden capas de Pb/Cu para $E>20$ MeV). Midiendo las tasas con todas las esferas
\[ N_d=\int R_d(E)\,\varphi(E)\,\mathrm{d}E \]
se obtiene el espectro $\varphi(E)$ por \textbf{deconvolución} (\emph{unfolding}, p.ej.\ código MAXED). Cubren de meV a GeV, pero con resolución pobre.
Ejemplo: NEMUS (PTB-Braunschweig).
\end{clave}
\begin{itemize}
\item \textbf{Long counter}: moderador + blindaje + BF$_3$; unidireccional, eficiencia casi plana con la energía.
\item \textbf{Matriz cilíndrica moderada} (contador $4\pi$): bloque de polietileno con muchos contadores de $^3$He/BF$_3$ alrededor de la fuente.
Alta eficiencia pero sin información de energía. Ej.: \textbf{NERO} (NSCL): $60\times60\times80$ cm$^3$, 16 contadores de $^3$He y 44 de BF$_3$,
$\varepsilon\approx40\%$ (emisión de neutrones $\beta$-retardados).
\item \textbf{Cámara de fisión}: depósito de $^{235}$U/$^{239}$Pu en una cámara de gas; los fragmentos ($\sim$160 MeV) dan pulsos enormes, ideal en campos $\gamma$ intensos.
\end{itemize}

\section{Discriminación y simulación}
\begin{itemize}
\item \textbf{Discriminación por altura de pulso}: umbrales para rechazar $\gamma$ y ruido ($^3$He excelente).
\item \textbf{PSD}: la forma del pulso depende de la partícula (BC501, NE213, CLYC, LaCl$_3$).
\item La respuesta se calcula por \textbf{Monte Carlo} (MCNP, Geant4, NRESP, SCINFUL), que necesita datos nucleares y la respuesta en luz del material.
El método MC moderno nació precisamente para el transporte de neutrones.
\end{itemize}

\setcounter{chapter}{6}
\chapter{Electrónica nuclear}
\begin{dato}
Este tema corresponde a las diapositivas ``Tema 9 -- Electrónica nuclear'' (J.M. Udías). Se coloca aquí porque completa
la cadena de detección (amplificación, conformado y digitalización) de los Temas 2--6.
\end{dato}

\section{Objetivos y cadena de lectura}
\begin{clave}[Cadena electrónica]
\begin{center}Detector (+HV) $\to$ Preamplificador $\to$ Amplificador/conformador $\to$ ADC/TDC, contador o digitalizador\end{center}
\begin{itemize}
\item La señal del detector es una \textbf{corriente}; su integral es la \textbf{carga}, proporcional a la energía depositada.
\item Objetivos: detectar pulsos individuales y medir su \textbf{energía} (amplitud/área), su \textbf{forma} y su \textbf{tiempo} de llegada (trigger, coincidencias, timing).
\item Funciones: adquisición de la señal, conformado, digitalización y lectura de datos.
\end{itemize}
\end{clave}
La mayoría de detectores en física nuclear opera en modo pulso; con flujos enormes, en modo corriente.

\section{Preamplificador}
\begin{clave}[Preamplificador sensible a carga]
Dos tipos: \textbf{sensible a tensión} ($V_o=-\frac{R_2}{R_1}V_i$, depende de la capacidad del detector) y
\textbf{sensible a carga} (el estándar para detectores de semiconductor):
\[ V_o\approx-\frac{Q}{C_f},\qquad V(t)=V_0\,e^{-t/R_fC_f} \]
La amplitud de salida sólo depende de la carga y de la capacidad de realimentación $C_f$, \emph{no} de la capacidad del
detector (que puede variar con la tensión o la temperatura). La resistencia $R_f$ descarga lentamente $C_f$ (cola de
$\sim$50--500 $\mu$s).
\begin{itemize}
\item Debe ser lineal (conservar la proporcionalidad carga--energía) y estar \textbf{lo más cerca posible del detector} (minimizar capacidad y ruido).
\end{itemize}
\end{clave}
Amplificador operacional ideal: $V_+-V_-=0$, sin corriente por las entradas; la ganancia la fija la realimentación (válido mientras la
frecuencia de la señal sea menor que el ancho de banda).

\section{Amplificador y conformado}
\begin{clave}[¿Por qué conformar?]
La salida del preamplificador tiene una cola larga ($\sim500\,\mu$s): a tasas $>100$ s$^{-1}$ los pulsos se montan unos sobre la cola de otros
(\textbf{apilamiento}) y la amplitud deja de medir la energía. El amplificador:
\begin{itemize}
\item \textbf{Acorta} los pulsos (reduce apilamiento).
\item \textbf{Mejora la relación señal/ruido} filtrando frecuencias: el \textbf{diferenciador CR} (paso alto, $\tau_d=C_dR_d$) elimina el ruido de baja
frecuencia; el \textbf{integrador RC} (paso bajo, $\tau_i=R_iC_i$) elimina el de alta frecuencia.
\item CR--RC$^n$ (varias integraciones) da un pulso \textbf{semigaussiano}; también conformado trapezoidal (digital).
\end{itemize}
\end{clave}
\begin{itemize}
\item \textbf{Déficit balístico}: si el tiempo de pico del conformador es menor que el tiempo de colección de carga, se pierde señal
(importante en HPGe grandes con tiempos de colección variables).
\item \textbf{Desplazamiento de línea base} (undershoot del pulso unipolar): se corrige con \textbf{cancelación polo-cero}.
\item \textbf{Apilamiento}: se reduce con constantes de tiempo cortas, a costa de peor S/N y resolución. Compromiso entre tasa y resolución.
\end{itemize}

\section{Ruido}
Ruido intrínseco: térmico (resistencias), de disparo (\emph{shot}, corriente de fuga), $1/f$. Se expresa como carga equivalente de ruido (ENC).
Depende de la capacidad del detector, la corriente de fuga y la constante de conformado (existe un $\tau$ óptimo).

\section{Medidas de tiempo}
\begin{clave}[Discriminadores y time walk]
\begin{itemize}
\item \textbf{Discriminador de umbral fijo} (\emph{leading edge}): marca el tiempo cuando la señal cruza un umbral.
Problema: \textbf{time walk}: pulsos de mayor amplitud cruzan antes que los pequeños (dependencia tiempo--energía).
\item \textbf{Discriminador de fracción constante (CFD)}: se suma la señal retrasada con una copia invertida y atenuada (fracción $f$);
el cruce por cero ocurre a una fracción fija de la amplitud $\Rightarrow$ \textbf{independiente de la amplitud} (elimina el walk para pulsos de igual forma).
\item \textbf{Jitter}: el ruido de tensión se traduce en incertidumbre temporal
\[ \Delta t\approx\frac{\Delta U}{U}\,t_r\qquad(\text{pendiente }U/t_r) \]
$\Rightarrow$ buen timing requiere \textbf{flancos rápidos} y alta S/N. Ej.: $U=100$ mV, ruido 1 mV, $t_r=1$ ns $\Rightarrow\Delta t\approx10$ ps.
\end{itemize}
\end{clave}
Un filtro paso bajo reduce el ruido manteniendo la mayor parte de la señal. Los TDC actuales tienen LSB de $\sim$10 ps.
Diferencias de tiempo: TAC (convertidor tiempo-amplitud) + ADC, o TDC.

\section{Análisis de alturas de pulso}
\begin{itemize}
\item \textbf{SCA} (analizador monocanal): discriminador inferior (LLD) y superior (ULD) en anticoincidencia; da un pulso lógico si la amplitud está
en la ventana $[E,E+\Delta E]$ (se pierde la información de energía). Alimenta contadores (\emph{scalers}).
\item \textbf{MCA} (multicanal): un \textbf{ADC} asigna cada pulso a un canal y se acumula un histograma = \textbf{espectro}.
\item \textbf{ADC}: \textbf{Wilkinson} (rampa: se carga un condensador y se cuenta el tiempo de descarga con un reloj; lineal, lento) y
\textbf{aproximaciones sucesivas} (rápido); también \emph{flash} (comparadores en paralelo). El ADC es el corazón del MCA y suele fijar el tiempo muerto.
\end{itemize}
\begin{dato}[Estándares]
\textbf{NIM} (módulos con alimentación $\pm6$, $\pm12$, $\pm24$ V; lógica NIM rápida), CAMAC (1 MB/s), \textbf{VME} (hasta 40 MB/s), y hoy digitalizadores en $\mu$TCA/PCIe.
\end{dato}

\section{Adquisición totalmente digital (FD-DAQ)}
\begin{clave}[Digitalización completa del pulso]
Se digitaliza la señal casi en bruto (ADC rápido, ``digitizer'' de cientos de MS/s a GS/s) y se procesa \textbf{numéricamente} para obtener energía y tiempo.
\begin{itemize}
\item \textbf{Ventajas}: sencillez (una tarjeta para energía y tiempo/coincidencias), conserva la forma del pulso (PSD offline),
flexibilidad (cualquier filtro: trapezoidal, mediana, FFT\dots), estabilidad e inmunidad al ruido una vez digitalizado.
\item \textbf{Inconvenientes}: otra cultura técnica, tecnología que evoluciona muy rápido, volumen de datos.
\item \textbf{Energía}: integración del pulso menos línea base, conformado trapezoidal o semigaussiano digital; corrección de déficit balístico trivial.
Resultados equivalentes a la electrónica analógica (LaBr$_3$: 4.6\% a 662 keV en ambos).
\item \textbf{Tiempo}: umbral fijo con interpolación (con walk), CFD digital, \textbf{umbrales relativos} (fracción de la amplitud; fácil en digital),
o algoritmos de optimización/\emph{machine learning} (mejora $\sim$15\% del CRT con LaBr$_3$).
\end{itemize}
\end{clave}
\begin{dato}
DRS4 (PSI): 4 canales, 1024 muestras, 1--5 GS/s, 12 bit. FATIMA (fast timing, UCM): LaBr$_3$(Ce) troncocónicos + PMT R9779.
Comparativa clásica: montaje de coincidencias analógico con preamplificador, amplificador, CFD, TAC y puerta frente a un único digitalizador.
\end{dato}

\setcounter{chapter}{8}
\chapter[Simulaciones Monte Carlo]{Introducción a las simulaciones\\ Monte Carlo}

\section{Qué es el método Monte Carlo}
\begin{clave}[Definición y componentes]
\textbf{Simulación estadística} de procesos físicos mediante el \textbf{muestreo de las variables aleatorias} que los caracterizan,
a partir de números aleatorios uniformes en $(0,1)$. Componentes:
\begin{enumerate}
\item \textbf{Funciones de densidad de probabilidad} que describen la física (secciones eficaces).
\item \textbf{Generador de números aleatorios} uniformes en $(0,1)$.
\item \textbf{Método de muestreo} de las variables a partir de esas distribuciones.
\item \textbf{Conteo} (\emph{scoring/tallies}): acumular las magnitudes de interés.
\item \textbf{Estimación del error} estadístico en función del número de historias.
\item Opcional: \textbf{reducción de varianza} y paralelización (el método es aditivo).
\end{enumerate}
\end{clave}
\begin{dato}[Historia]
Aguja de Buffon (1777): probabilidad de cruzar una línea $=2/\pi$. El método moderno nace en el Proyecto Manhattan
(Los Álamos; Ulam, von Neumann, Metropolis) para el transporte de neutrones.
\end{dato}

\section{Números aleatorios}
En un ordenador son \textbf{pseudoaleatorios}: los genera un algoritmo determinista a partir de una \textbf{semilla}; la misma
semilla da la misma secuencia (usar semillas distintas en cada simulación independiente). Requisitos:
\textbf{distribución uniforme}, \textbf{periodo largo} y \textbf{reproducibilidad}.

\section{Muestreo}
Se quiere generar $x\in(x_{min},x_{max})$ con densidad $p(x)\ge0$ normalizada, $\int p(x)\,\mathrm{d}x=1$.
\begin{clave}[Métodos de muestreo]
\textbf{Transformada inversa}: con $\xi$ uniforme en $(0,1)$,
\[ F(x)=\int_{x_{min}}^{x}p(x')\,\mathrm{d}x'=\xi\quad\Rightarrow\quad x=F^{-1}(\xi) \]
Ejemplo clave, \textbf{recorrido libre} hasta la siguiente interacción:
\[ p(s)=\mu e^{-\mu s}\ \Rightarrow\ s=-\frac{\ln\xi}{\mu}=-\lambda\ln\xi \]
\textbf{Rechazo} (cuando $F^{-1}$ no es analítica): se genera $x$ uniforme y $y$ uniforme en $(0,p_{max})$; se acepta $x$ si $y<p(x)$.
(Es el mismo principio que estimar $\pi$ o áreas lanzando puntos.)\\
\textbf{Variable discreta}: con probabilidades $P_i$, se divide $(0,1)$ en tramos de longitud $P_i$ y se elige el tramo donde cae $\xi$.
\end{clave}

\begin{ejemplo}[Elección de la interacción]
Fotón de 2 MeV: $\mu_C=0.1$, $\mu_{pp}=0.01$, $\mu_{fe}=0.0001$ cm$^{-1}$. Probabilidades $P_i=\mu_i/\mu_{tot}$:
Compton $[0,0.90)$, pares $[0.90,0.999)$, fotoeléctrico $[0.999,1)$. Con $\xi=0.418\to$ Compton; con $\xi=0.925\to$ pares.
\end{ejemplo}

\section{Incertidumbre estadística}
\begin{clave}[Error $\propto1/\sqrt N$]
Los resultados MC siempre llevan incertidumbre estadística, que decrece como
\[ \sigma\propto\frac{1}{\sqrt N}\qquad(N=\text{número de historias}) \]
Para reducirla a la mitad hay que simular \textbf{4 veces más} historias; para dividirla por 10, \textbf{100 veces más}.

\textbf{Eficiencia simulada} ($n$ de $N$ fotones en el fotopico): $\varepsilon=n/N$, distribución binomial:
\[ \sigma_\varepsilon=\sqrt{\frac{\varepsilon(1-\varepsilon)}{N}},\qquad \frac{\sigma_\varepsilon}{\varepsilon}=\sqrt{\frac{1-\varepsilon}{\varepsilon N}}\approx\frac{1}{\sqrt n}\ (\varepsilon\ll1) \]
\end{clave}
\begin{ejemplo}
$\gamma$ de 662 keV en NaI ($\mu=0.282$ cm$^{-1}$), $\xi=0.5\Rightarrow s=\ln2/0.282\approx2.5$ cm.
$N=10^6$, $n=2\cdot10^5$: $\varepsilon=0.2000\pm0.0004$ (0.2\%). Con $\varepsilon=2\%$ y $N=10^6$: $\sigma_\varepsilon/\varepsilon\approx0.7\%$; para 0.1\% harían falta $5\cdot10^7$ historias.
\end{ejemplo}

\paragraph{Ventajas e inconvenientes.} \emph{Pros}: problemas complejos sin simplificar; la precisión mejora simplemente simulando
más; no hace falta formular las ecuaciones de las magnitudes de interés. \emph{Contras}: alto coste computacional y resultados con incertidumbre estadística.

\section{Transporte de radiación}
\begin{clave}[Algoritmo básico (bucle sobre partículas primarias)]
\begin{enumerate}
\item Inicializar posición, dirección y energía de la partícula.
\item Si está en vacío, propagar hasta la primera interfase.
\item Calcular la sección eficaz total $\sigma$ (o $\mu$) a esa energía y material.
\item Recorrido libre medio $\lambda=1/(N\sigma)$.
\item Muestrear la distancia a la siguiente interacción: $p(s)=\frac1\lambda e^{-s/\lambda}\Rightarrow s=-\lambda\ln\xi$.
\item Elegir el proceso: $P_i=\sigma_i/\sigma$.
\item Muestrear pérdida de energía y nueva dirección de la sección eficaz diferencial del proceso $i$.
\item Añadir los secundarios a la pila (\emph{stack}).
\item Registrar la contribución a los observables (energía depositada, tiempo\dots).
\item Volver a 3 hasta que la energía baje del umbral (\emph{cut}) o la partícula salga del volumen de interés.
\end{enumerate}
\end{clave}
Para partículas cargadas, en lugar de simular cada colisión se usan \textbf{historias condensadas} (pasos con pérdida continua + dispersión múltiple) y
sólo se simulan explícitamente los sucesos ``duros'' (rayos $\delta$, bremsstrahlung por encima de un umbral).

\section{Reducción de varianza}
\begin{clave}[Idea]
Coste computacional $\propto\sigma^2\times t$ (varianza $\times$ tiempo de CPU por primaria). Las técnicas de \emph{biasing} reducen $\sigma$ o $t$
(a menudo una a costa de la otra); hay que \textbf{minimizar el producto}. MC \textbf{sesgado}: se muestrea de distribuciones artificiales y se
corrige con un \textbf{peso estadístico}, de modo que el valor esperado no cambia. Misma media con menos varianza, pero \emph{no} reproduce
correlaciones ni fluctuaciones suceso a suceso (el MC \textbf{analógico} sí).
\end{clave}
\begin{itemize}
\item \textbf{Forzado de interacciones / sesgo del camino libre}: aumentar artificialmente procesos poco probables (bremsstrahlung, fotonucleares, blancos finos).
\item \textbf{Particionado} (\emph{splitting}): al entrar en la región de interés una partícula se divide en $n$ con peso $w/n$ (reduce $\sigma$, aumenta $t$).
\item \textbf{Ruleta rusa}: al salir de la región de interés se elimina con cierta probabilidad, aumentando el peso de las supervivientes (reduce $t$, aumenta $\sigma$).
\item \textbf{Sesgo por importancia}: el usuario asigna importancias relativas a regiones (combina splitting y ruleta rusa). Útil en blindajes gruesos.
\item \textbf{Sesgo de la partícula principal} (cascadas EM): sólo se sigue uno de los dos secundarios, con probabilidad $\propto$ energía.
\end{itemize}

\section{Principales códigos}
\begin{dato}[Códigos de transporte]
\begin{tabularx}{\linewidth}{lX}
\textbf{Geant4} & CERN; todo tipo de partículas, 250 eV--10 TeV; geometrías complejas; C++; código abierto. Física nuclear, altas energías, física médica.\\
\textbf{MCNP} & Los Álamos; n, $\gamma$, $e^-$ (MCNPX/MCNP6: casi todas las partículas). \textbf{Referencia para neutrones}. Licencia controlada.\\
\textbf{FLUKA} & CERN/INFN; $\sim$60 partículas, hasta 1000 TeV, neutrones térmicos, iones. Aceleradores y protección radiológica.\\
\textbf{EGSnrc} & Sólo electromagnético (e$^\pm$, $\gamma$), 1 keV--10 GeV; BEAMnrc; muy usado en física médica (radioterapia).\\
\textbf{PENELOPE} & e$^\pm$, $\gamma$, 50 eV--1 GeV; geometría con cuádricas; física médica.\\
\textbf{ETRAN} & Primer código de transporte de e$^-$ y $\gamma$.\\
Específicos & NRESP, SCINFUL (respuesta de centelleadores a neutrones), SRIM/TRIM (iones).\\
\end{tabularx}
\end{dato}

\section{Aplicaciones}
\begin{itemize}
\item \textbf{Detectores}: funciones respuesta y eficiencias (Bonner, NaI, HPGe, neutrones), ángulos sólidos de fuentes extensas.
\item \textbf{Preparación y análisis de experimentos}: optimizar geometría, distancias y tamaños; entender el fondo. Requiere una descripción detallada de la geometría y del haz.
\item \textbf{Física médica}: dosis en radioterapia (los métodos semianalíticos tipo \emph{pencil beam} fallan en tejidos heterogéneos), protonterapia, imagen (PET, SPECT, TC).
\item \textbf{Aceleradores e instalaciones}: blindajes, absorbedores de haz, activación y tiempos de enfriamiento (protección radiológica).
\item Fuera de la física: epidemiología, finanzas, producción, transporte.
\end{itemize}

\chapter{Aceleradores de partículas}

\section{Idea general}
Un acelerador aumenta la energía de partículas cargadas mediante \textbf{campos eléctricos} y controla su trayectoria con
\textbf{campos magnéticos} (el tubo de rayos catódicos de un televisor, 30 kV, es el ejemplo más sencillo). Más energía permite
acceder a escalas más pequeñas.

\begin{clave}[¿Cómo dar muchos MeV a un ion sin muchos megavoltios?]
\begin{center}\small
\begin{tabularx}{\linewidth}{XXX}
\toprule
\textbf{Una sola vez} & \textbf{Muchas veces en línea} & \textbf{Muchas veces en círculo}\\
\midrule
Electrostáticos: $T=qV$. Tándem: el ion cambia de signo y aprovecha $V$ dos veces. &
Linac: cavidades de RF en fase con el haz. &
Ciclotrón: $B$ fijo, $f$ fija. Sincrotrón: $B$ y $f$ siguen al haz.\\
\midrule
Límite: \textbf{ruptura dieléctrica} & Límite: \textbf{longitud} & Límite: \textbf{relatividad} y \textbf{radiación sincrotrón}\\
\bottomrule
\end{tabularx}
\end{center}
\end{clave}

\section{Aceleradores electrostáticos}
Campos que no varían en el tiempo; la energía es proporcional a la tensión: $T=qV$.
\begin{itemize}
\item \textbf{Cockcroft--Walton}: multiplicador de tensión con diodos y condensadores (cada semiciclo carga el siguiente condensador).
Iones acelerados por una serie de electrodos a potenciales crecientes. En aire $\lesssim1$ MV (ruptura); en SF$_6$ a presión, varios MV.
Altas corrientes ($\sim$mA): también como inyector.
\item \textbf{Van de Graaff} (1929): una correa aislante transporta mecánicamente carga hasta una cúpula (la carga reside en la superficie
exterior del conductor). Varios MV en gas a presión.
\end{itemize}
\begin{clave}[Acelerador tándem]
Dos etapas con el terminal de alta tensión positiva $V$ en el centro:
\begin{enumerate}
\item Iones \textbf{negativos} (1 electrón extra) se aceleran desde tierra hasta el terminal: ganan $eV$.
\item En el terminal un \textbf{stripper} (lámina o gas) les arranca electrones: pasan a carga $+q$.
\item Se aceleran de nuevo del terminal a tierra: ganan $qV$.
\end{enumerate}
\[ T=(1+q)\,eV \]
Ejemplo: VdG de 10 MV: C$^-$ llega a 10 MeV y, como C$^{2+}$, sale con 30 MeV. Ventaja: fuente de iones y blanco a tierra.
\end{clave}
\begin{dato}
CMAM (Universidad Autónoma de Madrid): Tandetron de 5 MV con fuente Cockcroft--Walton. Terminales electrostáticos $\lesssim25$ MV.
\end{dato}

\section{Aceleradores de campo oscilante}
Las máquinas de tensión fija están limitadas por la ruptura; las \textbf{resonantes} reutilizan muchas veces un campo alterno moderado
(cavidades de radiofrecuencia, RF). Widerøe (1928) propuso aplicar una tensión alterna a una serie de huecos: concepto de linac.

\subsection{Acelerador lineal (linac)}
\begin{itemize}
\item Tubos metálicos en vacío conectados alternativamente a los terminales de un oscilador de RF. Dentro de los tubos no hay campo (jaula de Faraday);
la partícula se acelera en los huecos. El campo cambia de sentido mientras la partícula está dentro del tubo.
\item La \textbf{longitud de los tubos crece} con la velocidad para mantener la fase ($L_n=v_n/2f$). Cuando $v\to c$, cavidades equiespaciadas.
\item Variantes: DTL (tubos de deriva), RFQ (cuadrupolo de RF, primera etapa para iones lentos: acelera, agrupa y focaliza).
\item No hay límite teórico de energía (se vuelve a tierra tras cada cavidad); el límite es la \textbf{longitud}. Sin pérdidas por radiación: ideal para electrones de alta energía.
\end{itemize}
\begin{dato}
ILC (propuesto): linacs de $\sim$20 km para $e^\pm$ hasta 1 TeV. Linac4 (CERN). Linacs médicos de electrones (radioterapia, 6--25 MeV).
\end{dato}

\subsection{Ciclotrón}
\begin{clave}[Frecuencia ciclotrón]
Dos electrodos huecos en forma de \textbf{D} conectados a una tensión alterna, dentro de un campo magnético intenso uniforme. Sólo hay
campo eléctrico en el hueco entre las D. La fuerza de Lorentz da órbitas circulares:
\[ \frac{mv^2}{r}=qvB\ \Rightarrow\ r=\frac{mv}{qB}=\frac{p}{qB},\qquad f=\frac{qB}{2\pi m} \]
\textbf{La frecuencia no depende del radio ni de la velocidad} (no relativista): la partícula describe una espiral y llega siempre en fase al hueco.
Energía final (radio de extracción $R$):
\[ T=\frac{q^2B^2R^2}{2m} \]
Para protones con $B=1$ T, $f\approx15$ MHz (RF).
\end{clave}
\begin{dato}[Historia]
Lawrence (1929--31): primer ciclotrón de 4'' (11 cm), protones de 80 keV con sólo 2 kV; 1932, 11'' y 1.1 MeV. Nobel 1939.
TRIUMF (Vancouver): el mayor, 18 m, p (H$^-$) hasta 520 MeV. GANIL: 5 ciclotrones (C01, C02, CSS1, CSS2, SPIRAL1). Ciclotrones de 17--18 MeV para producir radioisótopos PET.
\end{dato}

\paragraph{Limitación relativista.} $\omega=qB/(\gamma m)$ disminuye con la energía: las partículas llegan tarde al hueco y se desfasan.
Soluciones:
\begin{itemize}
\item \textbf{Ciclotrón isócrono}: $B(r)$ crece con el radio $\propto\gamma(r)$ para mantener $f$ constante; haz continuo (CW). La desfocalización
vertical se compensa con campo variable azimutalmente (sectores, colinas/valles, espirales: focalización de Thomas). Ej.: PSI 590 MeV, IBA C230.
\item \textbf{Sincrociclotrón}: $B$ constante y $f$ variable: haz \textbf{pulsado}. Ej.: IBA Proteus One (protonterapia, 2.5 m, Quirónsalud Madrid).
\end{itemize}

\subsection{Betatrón}
Acelera \textbf{electrones} por \textbf{inducción}: un campo magnético variable (corriente alterna) produce un flujo variable que induce una fem
en una cámara toroidal; los electrones mantienen un radio constante. Kerst (1940): 2.3 MeV; hasta 300 MeV. Uso: rayos X de alta energía (blanco metálico).

\subsection{Sincrotrón}
\begin{clave}[Sincrotrón]
Para $v\approx c$:
\[ f=\frac{qB}{2\pi\gamma m},\qquad p=qBR \]
Radio \textbf{fijo}: imanes dipolares a lo largo del anillo y cavidades de RF pequeñas. Durante la aceleración \textbf{$B$ aumenta} con el momento
(y $f$ se ajusta para p/iones). Es el tipo de casi todos los aceleradores circulares modernos (LHC: 6.8 TeV por haz). Los colisionadores son sincrotrones.
\end{clave}
\begin{itemize}
\item \textbf{Radiación sincrotrón}: una carga relativista en órbita circular radia. Potencia $\propto\gamma^4/R^2$ (a igual energía, $\propto1/m^4$):
limita la energía de los sincrotrones de electrones (por eso para $e^\pm$ de muy alta energía se prefieren linacs) y, a la vez, es una potente
\textbf{fuente de rayos X} (ALBA, 3 GeV, Barcelona).
\item Necesita un inyector (linac o booster). Ejemplo: cadena del CERN (Linac4 $\to$ PSB $\to$ PS $\to$ SPS $\to$ LHC).
\end{itemize}

\section{Resumen}
\begin{clave}[Tipos de aceleradores]
\begin{center}\footnotesize
\begin{tabularx}{\linewidth}{>{\raggedright}p{2.6cm}>{\raggedright}p{2.1cm}>{\raggedright}p{2.0cm}>{\raggedright}X>{\raggedright\arraybackslash}p{2.4cm}}
\toprule
Tipo & Campo & Órbita/haz & Partículas, energía & Ejemplo\\
\midrule
Electrostático (C-W, VdG, tándem) & $E$ estático & Lineal, continuo & Iones; $\lesssim25$ MV & CMAM 5 MV\\
Linac de RF & $E$ de RF & Lineal & $e^-$ (MeV--GeV), p, iones & Linacs médicos, Linac4\\
Ciclotrón (clásico/isócrono) & $B$ cte, $f$ fija & Espiral, continuo & p, iones; $\lesssim600$ MeV & PET; PSI 590 MeV\\
Sincrociclotrón & $B$ cte, $f$ variable & Espiral, pulsado & p $\sim$250 MeV & IBA Proteus One\\
Betatrón & $E$ inducido ($B$ variable) & Radio fijo & $e^-\lesssim300$ MeV & Kerst (1940)\\
Sincrotrón & $B$ (y $f$) variable & Radio fijo & $e^-$ GeV, p hasta TeV & ALBA, LHC\\
\bottomrule
\end{tabularx}
\end{center}
\end{clave}

\paragraph{Aplicaciones médicas.} Protonterapia con sincrociclotrones (IBA Proteus One) o sincrotrones (Hitachi, Clínica Universidad de Navarra, Madrid);
ciclotrones para radioisótopos PET; linacs de electrones para radioterapia convencional.

\chapter{Medida de vidas medias en física nuclear}

\section{Ley de desintegración}
\begin{clave}[Ley exponencial]
La desintegración es un proceso \textbf{estocástico} a nivel de cada núcleo; el ritmo es proporcional al número de núcleos presentes:
\[
\frac{\mathrm{d}N}{\mathrm{d}t}=-\lambda N\ \Rightarrow\ N(t)=N_0e^{-\lambda t}=N_0e^{-t/\tau},\qquad
\tau=\frac1\lambda,\qquad T_{1/2}=\tau\ln2,\qquad A=\lambda N
\]
La vida media es una \textbf{propiedad de cada nivel} (estado fundamental o excitado). Anchura del nivel: $\Gamma=\hbar/\tau$
($\hbar=6.58\cdot10^{-16}$ eV s).
\end{clave}
\begin{itemize}
\item \textbf{Ramas de desintegración}: si un nivel decae por varios caminos, $\lambda=\sum_i\lambda_i$. \emph{Todas} las ramas muestran la misma
pendiente temporal $T_{1/2}$ (la del nivel). Vida media parcial y cociente de ramificación:
\[ BR_i=\frac{\lambda_i}{\lambda},\qquad T_{1/2}^{(i)}=\frac{T_{1/2}}{BR_i} \]
Para una transición $\gamma$ hay que tener en cuenta además la conversión interna: $\lambda_\gamma=\lambda_i/(1+\alpha)$.
\item \textbf{Vidas muy largas}: no se puede esperar a ver la caída; se mide la actividad $A=\lambda N$ conociendo $N$ (o se usa el
\textbf{equilibrio secular} en una cadena: $\lambda_1N_1=\lambda_2N_2$). Ej.: $^{230}$Th y $^{234}$U para datación (245\,250 y 75\,690 años).
\end{itemize}

\section{Elementos de matriz de transición}
\begin{clave}[¿Por qué medir vidas medias de estados excitados?]
La vida media de un estado que decae por emisión $\gamma$ da acceso directo a la \textbf{probabilidad de transición reducida} $B(\sigma L)$, que es un
elemento de matriz \emph{no diagonal} del operador electromagnético y está ligado a la \textbf{estructura}: carácter de partícula independiente o
\textbf{colectivo}, \textbf{forma y deformación} nuclear.
\[
\lambda(\sigma L)=\frac{8\pi(L+1)}{\hbar L[(2L+1)!!]^2}\left(\frac{E_\gamma}{\hbar c}\right)^{2L+1}B(\sigma L)
\]
\begin{itemize}
\item La tasa depende fuertemente de la energía ($E_\gamma^{2L+1}$) y de la multipolaridad $L$. Ejemplo: $\lambda(E2)[\mathrm{s^{-1}}]\approx1.22\cdot10^9\,E_\gamma^5[\mathrm{MeV}]\,B(E2)[e^2\mathrm{fm}^4]$.
\item Procedimiento: medir $T_{1/2}$ $\to$ vida parcial de la transición (BR, $\alpha$, mezcla multipolar) $\to$ quitar la dependencia energética: $B(\sigma L)$
$\to$ expresar en \textbf{unidades Weisskopf} (W.u.) para quitar la dependencia en $A$.
\end{itemize}
\end{clave}

\paragraph{Estimaciones de Weisskopf.} Estimación burda suponiendo que la transición la hace un solo nucleón entre capas
(funciones de onda de partícula independiente): $B_W(E2)=0.0594\,A^{4/3}\ e^2\mathrm{fm}^4$, $B_W(E1)=0.0645\,A^{2/3}\ e^2\mathrm{fm}^2$, $B_W(M1)=1.79\ \mu_N^2$.
\begin{itemize}
\item $B\sim1$ W.u.: transición de partícula independiente. $B\gg1$ W.u. (decenas--cientos en E2): \textbf{colectiva} (muchos nucleones coherentes).
$B\ll1$ W.u.: transición inhibida (isómeros, cambios de estructura).
\end{itemize}

\paragraph{Deformación.} Para núcleos deformados (parámetro $\beta$), el momento cuadrupolar intrínseco $Q_0$ se relaciona con $B(E2)$ de las transiciones
colectivas:
\[
B(E2;0^+\to2^+)=\frac{5}{16\pi}e^2Q_0^2,\qquad Q_0\approx\frac{3}{\sqrt{5\pi}}ZR_0^2\beta\,(1+0.16\beta)
\]
de modo que la vida media del primer $2^+$ de los pares--pares mide la deformación. Aplicaciones: evolución de capas, simetría especular,
sistemática de $B(E2)$, modos colectivos, coexistencia de formas.

\paragraph{Excitación coulombiana (COULEX).} Método complementario y directo: la excitación es puramente electromagnética, y da los elementos de
matriz (incluidos signos relativos y momentos cuadrupolares espectroscópicos $Q_s$: distingue prolato de oblato) a partir de energías e intensidades.

\section{Técnicas de medida}
\begin{clave}[Técnica según la escala de tiempo]
Idea general: medidas \textbf{start/stop} y ajuste de la distribución temporal $N(t)\propto e^{-t/\tau}$.
\begin{center}\small
\begin{tabular}{p{2.6cm}p{7.2cm}p{3.6cm}}
\toprule
& Técnica & Rango\\
\midrule
Directas (electrónicas) & Irradiación y contaje & $>$ minutos\\
& Espectroscopía etiquetada (implant.--desint.) & $>\mu$s\\
& Haz pulsado & ns -- ms\\
& Coincidencias $\gamma$--$\gamma$ & ns -- 100 ns\\
& \textbf{Fast timing} (centelleadores LaBr$_3$) & $>10$ ps\\
\midrule
Indirectas (Doppler) & \textbf{RDDS} (plunger) & ps -- ns\\
& \textbf{DSAM} & 100 fs -- ps\\
& FDS & 10 -- 100 fs\\
\midrule
Otras & COULEX, anchuras de resonancia ($\Gamma=\hbar/\tau$) & $<$ fs\\
\bottomrule
\end{tabular}
\end{center}
\end{clave}

\subsection{Técnicas electrónicas directas}
\begin{itemize}
\item \textbf{Irradiación y contaje} ($>1$ min): se irradia la muestra para producir el núcleo (fundamental o isómero) y se mide su actividad en un lugar de bajo fondo.
\item \textbf{Espectroscopía etiquetada por retroceso} ($>\mu$s): diferencia de tiempo entre la implantación de un núcleo (identificado en un separador) en un detector
de Si y su desintegración ($\alpha$, p, $\gamma$\dots).
\item \textbf{Haz pulsado} (ns--ms): el pulso del haz da el start, el $\gamma$ el stop (isómeros).
\item \textbf{Coincidencias retardadas}: start con una radiación que puebla el nivel, stop con la que lo despuebla.
\end{itemize}

\subsection{Métodos Doppler}
\begin{clave}[Principio]
Los núcleos se producen con gran velocidad de retroceso ($\beta\sim$ unos \%) en fusión--evaporación o COULEX en cinemática inversa.
Con $\beta=1\%$ recorren $1\,\mu$m en 3 ps: la \textbf{distancia recorrida es la escala de tiempo}. Los $\gamma$ emitidos en vuelo están desplazados:
\[ E_\gamma=E_0\left(1+\beta\cos\theta\right)\quad(\text{a primer orden}) \]
y se distinguen de los emitidos en reposo. \textbf{Sólo con HPGe} (hace falta gran resolución en energía).
\end{clave}
\begin{itemize}
\item \textbf{RDDS (plunger)}: blanco fino y, a distancia variable $d$ (motor lineal, $\mu$m--mm), un \emph{stopper} (Au) que para los núcleos. Se observan dos picos:
desplazado (emitido en vuelo, $I_s$) y no desplazado (emitido parado, $I_u$). La fracción que decae en vuelo da la vida media:
\[ \frac{I_u}{I_u+I_s}=e^{-d/(v\tau)} \]
Variante \textbf{diferencial (DDCM)}: se usa un degradador en vez de stopper (dos velocidades $\beta_1>\beta_2$) y coincidencias $\gamma$--$\gamma$
(``gating from above'' elimina la alimentación lateral). Se mejora el pico/fondo con detectores auxiliares (separador de retrocesos, detectores de neutrones).
\item \textbf{DSAM}: blanco sobre un soporte grueso de $Z$ alto: el núcleo se frena en el soporte mientras emite. La \textbf{forma de línea} (mezcla de
energías Doppler entre $\beta_0$ y 0) depende de la relación entre $\tau$ y el tiempo de frenado ($\sim$ps); se ajusta con simulaciones del frenado.
\item \textbf{FDS}: para $\tau<100$ fs (bandas superdeformadas) el núcleo emite antes de frenarse apreciablemente; se mide la fracción de desplazamiento
$F(\tau)=\langle\beta\rangle/\beta_0$ en función del ángulo.
\end{itemize}

\subsection{Fast timing}
\begin{clave}[Método de retardo temporal avanzado ($\beta\gamma\gamma$ / $\gamma\gamma\gamma$)]
Combina detectores con funciones complementarias:
\begin{itemize}
\item \textbf{HPGe}: alta resolución en energía para \textbf{seleccionar la cascada} (puerta), pero mala respuesta temporal.
\item \textbf{LaBr$_3$(Ce)} (o BaF$_2$): respuesta muy rápida para el \textbf{timing}, con resolución en energía moderada (suficiente para seleccionar transiciones).
\item Centelleador plástico $\beta$ (en desintegraciones): rápido, da el start.
\end{itemize}
Análisis de la distribución de diferencias de tiempo:
\begin{itemize}
\item $\tau\gtrsim50$ ps: \textbf{deconvolución de la pendiente} de la cola exponencial (pendiente $=1/\tau$).
\item $\tau\sim5$--70 ps: \textbf{desplazamiento del centroide} respecto a la respuesta prompt: $\Delta C=\tau$.
\end{itemize}
La clave es la \textbf{calibración temporal}: dependencia del tiempo prompt con la energía (walk) del fotopico y del Compton.
\end{clave}
\begin{dato}
Método de H. Mach et al.\ (NIM A 280, 1989). Ejemplo reciente: desintegración de $^{128}$Cd (Llanos-Expósito et al., PRC 111, 2025). Arrays FATIMA y similares.
\end{dato}

\chapter{Herramientas de análisis}
Tema práctico: herramientas que se usan en las prácticas. Basta con saber \textbf{qué hace cada una y cuándo usarla}.

\begin{clave}[Qué herramienta para qué]
\begin{center}\small
\begin{tabularx}{\linewidth}{lX}
\toprule
\textbf{SRIM} & Poder de frenado (electrónico y nuclear) y alcances de \textbf{iones} en materiales; tablas y gráficos (eV--GeV).\\
\textbf{TRIM} & Parte \textbf{Monte Carlo} de SRIM: simula trayectorias ion a ion.\\
\textbf{NNDC / NuDat3 / ENSDF} & Datos nucleares de referencia: carta de núclidos, niveles, $\gamma$, vidas medias, intensidades.\\
\textbf{jsROOT} & Visor web de ficheros ROOT (prácticas de neutrones y de vidas medias).\\
\textbf{Gnuplot} & Representación y ajustes sencillos (lineal, gaussiana + fondo lineal) de espectros.\\
\textbf{NIST} (ESTAR/PSTAR/ASTAR, XCOM) & Poder de frenado de $e^-$, p, $\alpha$ y coeficientes de atenuación de fotones (Tema 1).\\
\bottomrule
\end{tabularx}
\end{center}
\end{clave}

\section{SRIM y TRIM}
\textbf{SRIM} (\emph{The Stopping and Range of Ions in Matter}, \url{http://srim.org}) estima:
\begin{itemize}
\item pérdida de energía electrónica y nuclear de iones en distintos materiales (incluidos compuestos);
\item alcances medios y sus distribuciones (straggling), perfiles de implantación;
\item desplazamientos atómicos y daño por radiación.
\end{itemize}
Calcula el frenado electrónico de forma parecida a Bethe--Bloch pero con \textbf{correcciones semiempíricas} (carga efectiva, capas) válidas a baja energía,
y el frenado nuclear con el potencial universal \textbf{ZBL} (Ziegler--Biersack--Littmark). No sigue partículas individuales.

\textbf{TRIM} es el Monte Carlo del paquete: simula colisiones binarias y dispersión múltiple, y da la pérdida de energía por paso, la distribución angular,
las trayectorias 3D, los defectos y cascadas de desplazamiento, el \emph{sputtering} y la distribución de energía depositada.

\begin{ejemplo}[Uso típico]
Espesor de un detector $\Delta E$ o de una ventana/degradador; energía con que llega una partícula tras atravesar el blanco; profundidad de implantación de un haz radiactivo.
\end{ejemplo}

\section{Bases de datos nucleares}
\textbf{NNDC} (National Nuclear Data Center, Brookhaven): \textbf{NuDat3} (carta de núclidos interactiva, búsqueda de $\gamma$ y niveles) y \textbf{ENSDF}
(evaluaciones de estructura nuclear). Se usan para identificar picos de un espectro, conocer intensidades absolutas (calibración de eficiencia),
energías de fuentes de calibración y vidas medias.

\section{ROOT y Gnuplot}
\begin{itemize}
\item \textbf{ROOT} (CERN): entorno de análisis de datos (histogramas, árboles, ajustes). \textbf{jsROOT}: visor web de ficheros \texttt{.root} sin instalar nada.
\item \textbf{Gnuplot}: gratuito, sencillo; representar espectros (CassyLab, CoMPASS) y hacer ajustes. Se usa con macros en las prácticas computacionales
(Monte Carlo y procesado de pulsos).
\end{itemize}

\begin{dato}[Ajuste típico de un fotopico]
Gaussiana + fondo lineal:
$f(E)=A\exp\!\left[-\frac{(E-E_0)^2}{2\sigma^2}\right]+a+bE$.
Área neta $=\sqrt{2\pi}A\sigma$ (eficiencia de fotopico, Tema 2); FWHM $=2.355\sigma$ (resolución); centroide $E_0$ (calibración en energía, lineal canal--energía con fuentes conocidas).
\end{dato}

\end{document}
